% Generated mechanically from the frozen examples-defos.tex target.
% Do not edit this generated file; edit the bound locale source instead.

% OK, start here.
%

\setcounter{chapter}{92}
\chapter{形变问题}


\phantomsection
\label{examples-defos-section-phantom}


\section{引言}
\label{examples-defos-section-introduction}

\noindent
本章旨在具体研究《形式形变理论》《形变理论》和《余切复形》三章中
发展的一般理论的若干例子。

\medskip\noindent
Schlessinger 的论文 \cite{Sch} 第 3 节也讨论了一些例子。






\section{形变问题的例子}
\label{examples-defos-section-examples}

\noindent
此处应加入的内容如下：
\begin{enumerate}
\item 概形的形变：
\begin{enumerate}
\item Rim--Schlessinger 条件；
\item 切空间的计算；
\item 无穷小形变的计算；
\item 仿射超曲面的形变范畴。
\end{enumerate}
\item 层的形变（例如固定 $X/S$、$S$ 的一个有限型点 $s$，以及
$X_s$ 上的拟凝聚层 $\mathcal{F}_s$）；
\item 代数空间的形变（与概形的形变非常相似，或许甚至更容易？）；
\item 映射的形变（例如概形之间的态射；可以固定靶与源中的一个或两个）；
\item 此处还需补充。
\end{enumerate}





\section{一般纲要}
\label{examples-defos-section-general}

\noindent
本节列出讨论下面几个例子时所遵循的步骤。

\medskip\noindent
步骤 I。在每一节中，我们固定一个诺特环 $\Lambda$，并固定一个有限环同态
$\Lambda \to k$，其中 $k$ 是域。照例，令
$\mathcal{C}_\Lambda = \mathcal{C}_{\Lambda, k}$ 为基范畴，参见
《形式形变理论》定义 \ref{formal-defos-definition-CLambda}。

\medskip\noindent
步骤 II。在每一节中，我们定义一个在 $\mathcal{C}_\Lambda$ 上
以群胚为纤维余纤维化的范畴 $\mathcal{F}$。有时，我们转而考虑函子
$F : \mathcal{C}_\Lambda \to \textit{Sets}$。

\medskip\noindent
步骤 III。我们说明 $\mathcal{F}$ 在何种程度上满足《形式形变理论》第
\ref{formal-defos-section-RS-condition} 节所讨论的
Rim--Schlessinger 条件 (RS)。类似地，我们可以讨论 $\mathcal{F}$
在何种程度上满足 (S1) 与 (S2)，或 $F$ 在何种程度上满足相应的
Schlessinger 条件 (H1) 与 (H2)。参见《形式形变理论》第
\ref{formal-defos-section-schlessinger-conditions} 节。

\medskip\noindent
步骤 IV。设 $x_0$ 为 $\mathcal{F}(k)$ 的对象，换言之，它是
$\mathcal{F}$ 在 $k$ 上的对象。本章使用记号
$$
\Deformationcategory_{x_0} = \mathcal{F}_{x_0}
$$
表示《形式形变理论》注
\ref{formal-defos-remark-localize-cofibered-groupoid} 中构造的预形变范畴。
若 $\mathcal{F}$ 满足 (RS)，则 $\Deformationcategory_{x_0}$ 是形变范畴
（《形式形变理论》引理 \ref{formal-defos-lemma-localize-RS}），
并满足 (S1) 与 (S2)（《形式形变理论》引理
\ref{formal-defos-lemma-RS-implies-S1-S2}）。若满足 (S1) 与 (S2)，
则一个重要问题是切空间
$$
T\Deformationcategory_{x_0} = T_{x_0}\mathcal{F} = T\mathcal{F}_{x_0}
$$
是否有限维（参见《形式形变理论》注
\ref{formal-defos-remark-tangent-space-cofibered-groupoid} 与定义
\ref{formal-defos-definition-tangent-space}）。有限维性保证
$\Deformationcategory_{x_0}$ 具有通用形式对象（《形式形变理论》引理
\ref{formal-defos-lemma-versal-object-existence}）。

\medskip\noindent
步骤 V。若 $\mathcal{F}$ 通过步骤 IV，则下一个问题是 $k$-向量空间
$$
\text{Inf}(\Deformationcategory_{x_0}) = \text{Inf}_{x_0}(\mathcal{F})
$$
（即 $x_0$ 的无穷小自同构空间）是否有限维。若是，则
$\Deformationcategory_{x_0}$ 可由 $\mathcal{C}_\Lambda$ 上函子范畴中的
光滑 pro-可表示群胚给出表示，参见《形式形变理论》定理
\ref{formal-defos-theorem-presentation-deformation-groupoid}。




\section{有限投射模}
\label{examples-defos-section-finite-projective-modules}

\noindent
本节只是热身。有限投射模当然不应有任何“模空间”。

\begin{example}[有限投射模]
\label{examples-defos-example-finite-projective-modules}
如下定义范畴 $\mathcal{F}$：
\begin{enumerate}
\item 对象是二元组 $(A,M)$，其中 $A$ 是 $\mathcal{C}_\Lambda$ 的对象，
$M$ 是有限投射 $A$-模；
\item 态射 $(f,g):(B,N)\to(A,M)$ 由 $\mathcal{C}_\Lambda$ 中的态射
$f:B\to A$ 与映射 $g:N\to M$ 组成，其中该映射是 $f$-线性的，并诱导同构
$N\otimes_{B,f}A\cong M$。
\end{enumerate}
函子 $p:\mathcal{F}\to\mathcal{C}_\Lambda$ 把 $(A,M)$ 送到 $A$，
把 $(f,g)$ 送到 $f$。显然，$p$ 以群胚为纤维余纤维化。
给定有限维 $k$-向量空间 $V$，令 $x_0=(k,V)$ 为
$\mathcal{F}(k)$ 的相应对象。记
$$
\Deformationcategory_V = \mathcal{F}_{x_0}
$$
\end{example}

\noindent
由于局部环上的每个有限投射模都是有限自由模（《代数》引理
\ref{algebra-lemma-finite-projective}），可见
$$
\begin{matrix}
\mathcal{F}(A)\text{ 的对象的} \\
\text{同构类}
\end{matrix}
= \coprod\nolimits_{n \geq 0} \{*\}
$$
尽管这意味着 $\mathcal{F}$ 的形变理论本质上是平凡的，我们仍依照
第 \ref{examples-defos-section-general} 节所列步骤完整推导，作为一个简单例子。

\begin{lemma}
\label{examples-defos-lemma-finite-projective-modules-RS}
例 \ref{examples-defos-example-finite-projective-modules} 满足 Rim--Schlessinger 条件
(RS)。特别地，对任意有限维 $k$-向量空间 $V$，
$\Deformationcategory_V$ 是形变范畴。
\end{lemma}

\begin{proof}
设 $A_1\to A$ 与 $A_2\to A$ 是 $\mathcal{C}_\Lambda$ 的态射，
并假设 $A_2\to A$ 满射。由《形式形变理论》引理
\ref{formal-defos-lemma-RS-2-categorical}，只须证明函子
$\mathcal{F}(A_1 \times_A A_2) \to
\mathcal{F}(A_1) \times_{\mathcal{F}(A)} \mathcal{F}(A_2)$
是范畴等价。

\medskip\noindent
因此，须证明 $A_1\times_AA_2$ 上的有限投射模范畴，等价于
$A_1$ 与 $A_2$ 上的有限投射模范畴在 $A$ 上有限投射模范畴之上的
纤维积。这是《更多代数》引理
\ref{more-algebra-lemma-finitely-presented-module-over-fibre-product}
的特殊情形。回忆，其逆函子把三元组 $(M_1,M_2,\varphi)$ 送到有限投射
$A_1\times_AA_2$-模 $M_1\times_\varphi M_2$；这里 $M_1$ 是有限投射
$A_1$-模，$M_2$ 是有限投射 $A_2$-模，而
$\varphi:M_1\otimes_{A_1}A\to M_2\otimes_{A_2}A$ 是 $A$-模同构。
\end{proof}

\begin{lemma}
\label{examples-defos-lemma-finite-projective-modules-TI}
在例 \ref{examples-defos-example-finite-projective-modules} 中，设 $V$ 为有限维
$k$-向量空间。则
$$
T\Deformationcategory_V = (0)
\quad\text{且}\quad
\text{Inf}(\Deformationcategory_V) = \text{End}_k(V)
$$
均为有限维。
\end{lemma}

\begin{proof}
令 $\mathcal{F}$ 如例 \ref{examples-defos-example-finite-projective-modules}，并置
$x_0=(k,V)\in\Ob(\mathcal{F}(k))$。回忆，
$T\Deformationcategory_V=T_{x_0}\mathcal{F}$ 是二元组 $(x,\alpha)$
的同构类之集；其中 $x$ 是 $\mathcal{F}$ 在对偶数 $k[\epsilon]$ 上的对象，
$\alpha:x\to x_0$ 是 $\mathcal{F}$ 中位于 $k[\epsilon]\to k$ 上方的态射。

\medskip\noindent
在同构意义下，存在唯一的二元组 $(M,\alpha)$，其中 $M$ 是
$k[\epsilon]$ 上的有限投射模，$\alpha:M\to V$ 是诱导同构
$M\otimes_{k[\epsilon]}k\to V$ 的 $k[\epsilon]$-线性映射。
例如，若 $V=k^{\oplus n}$，则取 $M=k[\epsilon]^{\oplus n}$，
并令 $\alpha$ 为显然的映射。

\medskip\noindent
类似地，$\text{Inf}(\Deformationcategory_V)=\text{Inf}_{x_0}(\mathcal{F})$
是 $x_0$ 在 $k[\epsilon]$ 上的平凡形变 $x'_0$ 的自同构之集。
详见《形式形变理论》定义
\ref{formal-defos-definition-infinitesimal-auts}。

\medskip\noindent
给定第二段中的 $(M,\alpha)$，可见 $\text{Inf}_{x_0}(\mathcal{F})$
的元素是满足 $\gamma\bmod\epsilon=\text{id}$ 的自同构
$\gamma:M\to M$。于是可写成
$\gamma=\text{id}_M+\epsilon\psi$，其中
$\psi:M/\epsilon M\to M/\epsilon M$ 是 $k$-线性的。利用 $\alpha$，
可将 $\psi$ 看作 $\text{End}_k(V)$ 的元素，证明完毕。
\end{proof}


\section{群表示}
\label{examples-defos-section-representations}

\noindent
表示的形变理论可能非常有趣。

\begin{example}[群表示]
\label{examples-defos-example-representations}
设 $\Gamma$ 为群。如下定义范畴 $\mathcal{F}$：
\begin{enumerate}
\item 对象是三元组 $(A,M,\rho)$，其中 $A$ 是 $\mathcal{C}_\Lambda$
的对象，$M$ 是有限投射 $A$-模，且
$\rho:\Gamma\to\text{GL}_A(M)$ 是同态；
\item 态射 $(f,g):(B,N,\tau)\to(A,M,\rho)$ 由
$\mathcal{C}_\Lambda$ 中的态射 $f:B\to A$ 与映射 $g:N\to M$
组成，其中该映射是 $f$-线性且 $\Gamma$-等变的，并诱导同构
$N\otimes_{B,f}A\cong M$。
\end{enumerate}
函子 $p:\mathcal{F}\to\mathcal{C}_\Lambda$ 把 $(A,M,\rho)$ 送到 $A$，
把 $(f,g)$ 送到 $f$。显然，$p$ 以群胚为纤维余纤维化。
给定有限维 $k$-向量空间 $V$ 与表示
$\rho_0:\Gamma\to\text{GL}_k(V)$，令 $x_0=(k,V,\rho_0)$ 为
$\mathcal{F}(k)$ 的相应对象。记
$$
\Deformationcategory_{V, \rho_0} = \mathcal{F}_{x_0}
$$
\end{example}

\noindent
由于局部环上的每个有限投射模都是有限自由模（《代数》引理
\ref{algebra-lemma-finite-projective}），可见
$$
\begin{matrix}
\mathcal{F}(A)\text{ 的对象的} \\
\text{同构类}
\end{matrix}
=
\coprod\nolimits_{n \geq 0}\quad
\begin{matrix}
\text{同态 }\rho:\Gamma\to\text{GL}_n(A)\text{ 的}\\
\text{GL}_n(A)\text{-共轭类}
\end{matrix}
$$
这已经比第 \ref{examples-defos-section-finite-projective-modules} 节的讨论更有趣。

\begin{lemma}
\label{examples-defos-lemma-representations-RS}
例 \ref{examples-defos-example-representations} 满足 Rim--Schlessinger 条件 (RS)。
特别地，对任意有限维表示 $\rho_0:\Gamma\to\text{GL}_k(V)$，
$\Deformationcategory_{V,\rho_0}$ 是形变范畴。
\end{lemma}

\begin{proof}
设 $A_1\to A$ 与 $A_2\to A$ 是 $\mathcal{C}_\Lambda$ 的态射，
并假设 $A_2\to A$ 满射。由《形式形变理论》引理
\ref{formal-defos-lemma-RS-2-categorical}，只须证明函子
$\mathcal{F}(A_1 \times_A A_2) \to
\mathcal{F}(A_1) \times_{\mathcal{F}(A)} \mathcal{F}(A_2)$
是范畴等价。

\medskip\noindent
考虑范畴 $\mathcal{F}(A_1)\times_{\mathcal{F}(A)}\mathcal{F}(A_2)$
的对象
$$
((A_1, M_1, \rho_1), (A_2, M_2, \rho_2), (\text{id}_A, \varphi))
$$
。如引理 \ref{examples-defos-lemma-finite-projective-modules-RS} 的证明所示，
可以考虑有限投射 $A_1\times_AA_2$-模 $M_1\times_\varphi M_2$。
由于 $\varphi$ 与给定作用相容，得到
$$
\rho_1 \times \rho_2 : \Gamma \longrightarrow
\text{GL}_{A_1 \times_A A_2}(M_1 \times_\varphi M_2)
$$
于是 $(M_1\times_\varphi M_2,\rho_1\times\rho_2)$ 是
$\mathcal{F}(A_1\times_AA_2)$ 的对象。此构造给出上述函子的拟逆。
\end{proof}

\begin{lemma}
\label{examples-defos-lemma-representations-TI}
在例 \ref{examples-defos-example-representations} 中，设
$\rho_0:\Gamma\to\text{GL}_k(V)$ 为有限维表示。则
$$
T\Deformationcategory_{V, \rho_0} = \Ext^1_{k[\Gamma]}(V, V) =
H^1(\Gamma, \text{End}_k(V))
\quad\text{且}\quad
\text{Inf}(\Deformationcategory_{V, \rho_0}) = H^0(\Gamma, \text{End}_k(V))
$$
因此，$\text{Inf}(\Deformationcategory_{V,\rho_0})$ 总是有限维；
若 $\Gamma$ 有限生成，则 $T\Deformationcategory_{V,\rho_0}$ 也是有限维。
\end{lemma}

\begin{proof}
先处理无穷小自同构。令 $M=V\otimes_k k[\epsilon]$，并赋予其诱导作用
$\rho_0':\Gamma\to\text{GL}_n(M)$。无穷小自同构，即
$\text{Inf}(\Deformationcategory_{V,\rho_0})$ 的元素，由自同构
$\gamma = \text{id} + \epsilon \psi : M \to M$
给出，正如引理 \ref{examples-defos-lemma-finite-projective-modules-TI} 的证明；
此外，$\psi$ 必须与 $\Gamma$ 的作用（由 $\rho_0$ 给出）交换。因此
$$
\text{Inf}(\Deformationcategory_{V, \rho_0}) = H^0(\Gamma, \text{End}_k(V))
$$
，正如引理所述。

\medskip\noindent
其次，设 $(k[\epsilon],M,\rho)$ 为 $\mathcal{F}$ 在 $k[\epsilon]$
上的对象，并设 $\alpha:M\to V$ 为诱导同构 $M/\epsilon M\to V$
的 $\Gamma$-等变映射。由于 $M$ 是自由 $k[\epsilon]$-模，得到
$\Gamma$-模的扩张
$$
0 \to V \to M \xrightarrow{\alpha} V \to 0
$$
左端映射的具体构造从略。反之，若给定如上的 $\Gamma$-模扩张，
便可据此在 $M$ 上赋予 $k[\epsilon]$-模结构，并得到
$\mathcal{F}(k[\epsilon])$ 的对象以及如上的映射 $\alpha$。因此
$$
T\Deformationcategory_{V, \rho_0} = \Ext^1_{k[\Gamma]}(V, V)
$$
，正如引理所述。由《Étale 上同调》引理
\ref{etale-cohomology-lemma-ext-modules-hom}，它等于
$H^1(\Gamma,\text{End}_k(V))$。

\medskip\noindent
关于维数的断言来自《Étale 上同调》引理
\ref{etale-cohomology-lemma-finite-dim-group-cohomology}。
\end{proof}

\noindent
在例 \ref{examples-defos-example-representations} 中，若 $\Gamma$ 有限生成，且
$(V,\rho_0)$ 是 $\Gamma$ 在 $k$ 上的有限维表示，则
$\Deformationcategory_{V,\rho_0}$ 可由 $\mathcal{C}_\Lambda$ 上
函子范畴中的光滑 pro-可表示群胚给出表示，因而具有一个（极小）通用
形式对象。这来自引理 \ref{examples-defos-lemma-representations-RS}、
\ref{examples-defos-lemma-representations-TI} 以及第 \ref{examples-defos-section-general} 节的一般讨论。

\begin{lemma}
\label{examples-defos-lemma-representations-hull}
在例 \ref{examples-defos-example-representations} 中，假设 $\Gamma$ 有限生成。
设 $\rho_0:\Gamma\to\text{GL}_k(V)$ 为有限维表示，并假设
$\Lambda$ 是剩余域为 $k$ 的完备局部环（经典情形）。则函子
$$
F : \mathcal{C}_\Lambda \longrightarrow \textit{Sets},\quad
A \longmapsto \Ob(\Deformationcategory_{V, \rho_0}(A))/\cong
$$
（它取对象的同构类）具有包络。若
$H^0(\Gamma,\text{End}_k(V))=k$，则 $F$ 是 pro-可表示的。
\end{lemma}

\begin{proof}
包络的存在性来自引理 \ref{examples-defos-lemma-representations-RS}、
\ref{examples-defos-lemma-representations-TI}、《形式形变理论》引理
\ref{formal-defos-lemma-RS-implies-S1-S2} 与注
\ref{formal-defos-remark-compose-minimal-into-iso-classes}。

\medskip\noindent
假设 $H^0(\Gamma,\text{End}_k(V))=k$。为证明 $F$ 是 pro-可表示的，
只须证明 $F$ 是形变函子，参见《形式形变理论》定理
\ref{formal-defos-theorem-Schlessinger-prorepresentability}。
换言之，须证明 $F$ 满足 (RS)。为此可使用《形式形变理论》引理
\ref{formal-defos-lemma-RS-associated-functor} 的判据。
若能证明
$$
A \cdot \text{id}_M =
\text{End}_{A[\Gamma]}(M)
$$
对 $\mathcal{F}$ 中每个满足 $M\otimes_Ak$ 作为 $\Gamma$-表示同构于
$V$ 的对象 $(A,M,\rho)$ 都成立，便可得到所需的自同构群满射性。
由于左端包含于右端，只须证明
$\text{length}_A \text{End}_{A[\Gamma]}(M) \leq \text{length}_A A$.
选取两两不同的理想
$(0) = I_n \subset \ldots \subset I_1 \subset A$
，其中 $n=\text{length}(A)$。相应地过滤 $M$，可见只须证明
$\Hom_{A[\Gamma]}(M,I_tM/I_{t+1}M)$ 的长度为 $1$。由于
$I_tM/I_{t+1}M\cong M\otimes_Ak$，且任意 $A[\Gamma]$-模映射
$M\to M\otimes_Ak$ 都唯一地通过商映射 $M\to M\otimes_Ak$ 分解，
从而给出
$$
\text{End}_{A[\Gamma]}(M \otimes_A k) = \text{End}_{k[\Gamma]}(V) = k
$$
，结论即得。
\end{proof}



\section{连续表示}
\label{examples-defos-section-continuous-representations}

\noindent
一个非常有趣的做法是选取无限 Galois 群，研究其表示的形变理论，
参见 \cite{Mazur-deforming}。

\begin{example}[拓扑群的表示]
\label{examples-defos-example-continuous-representations}
设 $\Gamma$ 为拓扑群。如下定义范畴 $\mathcal{F}$：
\begin{enumerate}
\item 对象是三元组 $(A,M,\rho)$，其中 $A$ 是 $\mathcal{C}_\Lambda$
的对象，$M$ 是有限投射 $A$-模，且
$\rho:\Gamma\to\text{GL}_A(M)$ 是连续同态；这里
$\text{GL}_A(M)$ 赋予离散拓扑\footnote{另一种做法是要求由 $\rho$
赋予 $G$-作用的 $A$-模 $M$，是《Étale 上同调》定义
\ref{etale-cohomology-definition-G-module-continuous} 所定义的
$A\text{-}G$-模。不过，由于 $M$ 是有限 $A$-模，两种要求等价。}；
\item 态射 $(f,g):(B,N,\tau)\to(A,M,\rho)$ 由
$\mathcal{C}_\Lambda$ 中的态射 $f:B\to A$ 与映射 $g:N\to M$
组成，其中该映射是 $f$-线性且 $\Gamma$-等变的，并诱导同构
$N\otimes_{B,f}A\cong M$。
\end{enumerate}
函子 $p:\mathcal{F}\to\mathcal{C}_\Lambda$ 把 $(A,M,\rho)$ 送到 $A$，
把 $(f,g)$ 送到 $f$。显然，$p$ 以群胚为纤维余纤维化。
给定有限维 $k$-向量空间 $V$ 与连续表示
$\rho_0:\Gamma\to\text{GL}_k(V)$，令 $x_0=(k,V,\rho_0)$ 为
$\mathcal{F}(k)$ 的相应对象。记
$$
\Deformationcategory_{V, \rho_0} = \mathcal{F}_{x_0}
$$
\end{example}

\noindent
由于局部环上的每个有限投射模都是有限自由模（《代数》引理
\ref{algebra-lemma-finite-projective}），可见
$$
\begin{matrix}
\mathcal{F}(A)\text{ 的对象的} \\
\text{同构类}
\end{matrix}
=
\coprod\nolimits_{n \geq 0}\quad
\begin{matrix}
\text{连续同态 }\rho:\Gamma\to\text{GL}_n(A)\text{ 的}\\
\text{GL}_n(A)\text{-共轭类}
\end{matrix}
$$

\begin{lemma}
\label{examples-defos-lemma-continuous-representations-RS}
例 \ref{examples-defos-example-continuous-representations} 满足 Rim--Schlessinger 条件
(RS)。特别地，对任意有限维连续表示
$\rho_0:\Gamma\to\text{GL}_k(V)$，
$\Deformationcategory_{V,\rho_0}$ 是形变范畴。
\end{lemma}

\begin{proof}
证明与引理 \ref{examples-defos-lemma-representations-RS} 的证明完全相同。
\end{proof}

\begin{lemma}
\label{examples-defos-lemma-continuous-representations-TI}
在例 \ref{examples-defos-example-continuous-representations} 中，设
$\rho_0:\Gamma\to\text{GL}_k(V)$ 为有限维连续表示。则
$$
T\Deformationcategory_{V, \rho_0} = H^1(\Gamma, \text{End}_k(V))
\quad\text{且}\quad
\text{Inf}(\Deformationcategory_{V, \rho_0}) = H^0(\Gamma, \text{End}_k(V))
$$
因此，$\text{Inf}(\Deformationcategory_{V,\rho_0})$ 总是有限维；
若 $\Gamma$ 拓扑有限生成，则 $T\Deformationcategory_{V,\rho_0}$
也是有限维。
\end{lemma}

\begin{proof}
证明与引理 \ref{examples-defos-lemma-representations-TI} 的证明完全相同。
\end{proof}

\noindent
在例 \ref{examples-defos-example-continuous-representations} 中，若 $\Gamma$ 拓扑有限生成，
且 $(V,\rho_0)$ 是 $\Gamma$ 在 $k$ 上的有限维连续表示，则
$\Deformationcategory_{V,\rho_0}$ 可由 $\mathcal{C}_\Lambda$ 上
函子范畴中的光滑 pro-可表示群胚给出表示，因而具有一个（极小）通用
形式对象。这来自引理 \ref{examples-defos-lemma-continuous-representations-RS}、
\ref{examples-defos-lemma-continuous-representations-TI} 以及第
\ref{examples-defos-section-general} 节的一般讨论。

\begin{lemma}
\label{examples-defos-lemma-continuous-representations-hull}
在例 \ref{examples-defos-example-continuous-representations} 中，假设 $\Gamma$
拓扑有限生成。设 $\rho_0:\Gamma\to\text{GL}_k(V)$ 为有限维表示，
并假设 $\Lambda$ 是剩余域为 $k$ 的完备局部环（经典情形）。则函子
$$
F : \mathcal{C}_\Lambda \longrightarrow \textit{Sets},\quad
A \longmapsto \Ob(\Deformationcategory_{V, \rho_0}(A))/\cong
$$
（它取对象的同构类）具有包络。若
$H^0(\Gamma,\text{End}_k(V))=k$，则 $F$ 是 pro-可表示的。
\end{lemma}

\begin{proof}
证明与引理 \ref{examples-defos-lemma-representations-hull} 的证明完全相同。
\end{proof}



\section{分次代数}
\label{examples-defos-section-graded-algebras}

\noindent
在证明极化固有概形之叠是代数叠时，我们将使用本节的例子。因此，
我们考虑齐次分量均为有限投射模的交换分次代数（有时称为“局部有限”）。

\begin{example}[分次代数]
\label{examples-defos-example-graded-algebras}
如下定义范畴 $\mathcal{F}$：
\begin{enumerate}
\item 对象是二元组 $(A,P)$，其中 $A$ 是 $\mathcal{C}_\Lambda$ 的对象，
$P$ 是分次 $A$-代数，并且对所有 $d\geq0$，$P_d$ 都是有限投射 $A$-模；
\item 态射 $(f,g):(B,Q)\to(A,P)$ 由 $\mathcal{C}_\Lambda$ 中的态射
$f:B\to A$ 与映射 $g:Q\to P$ 组成，其中该映射是 $f$-线性的，
并诱导同构 $Q\otimes_{B,f}A\cong P$。
\end{enumerate}
函子 $p:\mathcal{F}\to\mathcal{C}_\Lambda$ 把 $(A,P)$ 送到 $A$，
把 $(f,g)$ 送到 $f$。显然，$p$ 以群胚为纤维余纤维化。
给定满足对所有 $d\geq0$ 均有 $\dim_k(P_d)<\infty$ 的分次
$k$-代数 $P$，令 $x_0=(k,P)$ 为 $\mathcal{F}(k)$ 的相应对象。记
$$
\Deformationcategory_P = \mathcal{F}_{x_0}
$$
\end{example}

\begin{lemma}
\label{examples-defos-lemma-graded-algebras-RS}
例 \ref{examples-defos-example-graded-algebras} 满足 Rim--Schlessinger 条件 (RS)。
特别地，对任意分次 $k$-代数 $P$，$\Deformationcategory_P$ 是形变范畴。
\end{lemma}

\begin{proof}
设 $A_1\to A$ 与 $A_2\to A$ 是 $\mathcal{C}_\Lambda$ 的态射，
并假设 $A_2\to A$ 满射。由《形式形变理论》引理
\ref{formal-defos-lemma-RS-2-categorical}，只须证明函子
$\mathcal{F}(A_1 \times_A A_2) \to
\mathcal{F}(A_1) \times_{\mathcal{F}(A)} \mathcal{F}(A_2)$
是范畴等价。

\medskip\noindent
考虑范畴 $\mathcal{F}(A_1)\times_{\mathcal{F}(A)}\mathcal{F}(A_2)$
的对象
$$
((A_1, P_1), (A_2, P_2), (\text{id}_A, \varphi))
$$
。考虑 $P_1\times_\varphi P_2$。由于
$\varphi : P_1 \otimes_{A_1} A \to P_2 \otimes_{A_2} A$
是分次代数同构，可见 $P_1\times_\varphi P_2$ 的各分次分量都是有限投射
$A_1\times_AA_2$-模，参见引理
\ref{examples-defos-lemma-finite-projective-modules-RS} 的证明。因此，
$P_1\times_\varphi P_2$ 是 $\mathcal{F}(A_1\times_AA_2)$ 的对象。
此构造给出上述函子的拟逆，证明完毕。
\end{proof}

\begin{lemma}
\label{examples-defos-lemma-graded-algebras-TI}
在例 \ref{examples-defos-example-graded-algebras} 中，设 $P$ 为分次 $k$-代数。则
$$
T\Deformationcategory_P
\quad\text{且}\quad
\text{Inf}(\Deformationcategory_P) = \text{Der}_k(P, P)
$$
若 $P$ 在 $k$ 上有限生成，二者均为有限维。
\end{lemma}

\begin{proof}
先处理无穷小自同构。令 $Q=P\otimes_k k[\epsilon]$。
$\text{Inf}(\Deformationcategory_P)$ 的元素由自同构
$\gamma = \text{id} + \epsilon \delta : Q \to Q$
给出，其中 $\delta:P\to P$。$\gamma$ 是分次的，蕴含 $\delta$
是次数为 $0$ 的齐次映射；$\gamma$ 是 $k$-线性的，蕴含 $\delta$
是 $k$-线性的；$\gamma$ 保乘法，蕴含 $\delta$ 是 $k$-导子。
反之，给定次数为 $0$ 的齐次 $k$-导子 $\delta:P\to P$，便得到如上的
自同构 $\gamma=\text{id}+\epsilon\delta$。因此
$$
\text{Inf}(\Deformationcategory_P) = \text{Der}_k(P, P)
$$
，正如引理所述。显然，若 $P$ 由 $P_i$（$0\leq i\leq N$）中的元素生成，
则 $\delta$ 由线性映射 $\delta_i:P_i\to P_i$（$0\leq i\leq N$）确定，
因而
$$
\dim_k \text{Der}_k(P, P) < \infty
$$
，如所需。

\medskip\noindent
为完成引理的证明，我们证明形变空间是有限维的。为此选取分次
$k$-代数的呈示
$$
k[X_1, \ldots, X_n]/(F_1, \ldots, F_m) \longrightarrow P
$$
，其中 $\deg(X_i)=d_i$，且 $F_j$ 是次数为 $e_j$ 的齐次元。
设 $Q$ 为每一次数上有限自由的任意分次 $k[\epsilon]$-代数，并给定同构
$\alpha:Q/\epsilon Q\to P$，使 $(Q,\alpha)$ 定义
$T\Deformationcategory_P$ 的元素。选取次数为 $d_i$ 的齐次元
$q_i\in Q$，使其映到 $X_i$ 在 $P$ 中的像。于是得到
$$
k[\epsilon][X_1, \ldots, X_n] \longrightarrow Q,\quad
X_i \longmapsto q_i
$$
。由于 $P=Q/\epsilon Q$，Nakayama 引理表明此映射满射。
一个简短的追图表明，可以选取次数为 $e_j$ 的齐次元
$F_{\epsilon,j}\in k[\epsilon][X_1,\ldots,X_n]$，使其在 $Q$ 中映到零，
并在 $k[X_1,\ldots,X_n]$ 中映到 $F_j$。于是
$$
k[\epsilon][X_1, \ldots, X_n]/(F_{\epsilon, 1}, \ldots, F_{\epsilon, m})
\longrightarrow Q
$$
由 $Q$ 在 $k[\epsilon]$ 上的平坦性，这是 $Q$ 的一个呈示。写成
$$
F_{\epsilon, j} =  F_j + \epsilon G_j
$$
向量 $(G_1,\ldots,G_m)$ 存在某些不定性。首先，改变
$F_{\epsilon,j}$ 的选择，可以用核
映射 $k[X_1,\ldots,X_n]\to P$ 的核中次数为 $e_j$ 的任意元素修改 $G_j$。
因此，我们不记录 $(G_1,\ldots,G_m)$ 本身，而记录元素
$$
(g_1, \ldots, g_m) \in P_{e_1} \oplus \ldots \oplus P_{e_m}
$$
，其中 $g_j$ 是 $G_j$ 在 $P_{e_j}$ 中的像。此外，若把 $q_i$ 的选择改为
$q_i+\epsilon p_i$，其中 $p_i$ 的次数为 $d_i$，则一个从略的计算表明，
$g_j$ 变为
$$
g_j^{new} = g_j - \sum\nolimits_{i = 1}^n p_i \partial F_j / \partial X_i
$$
由此可见，$Q$ 的同构类由向量 $(G_1,\ldots,G_m)$ 在 $k$-向量空间
$$
W  = \Coker(P_{d_1} \oplus \ldots \oplus P_{d_n}
\xrightarrow{(\frac{\partial F_j}{\partial X_i})}
P_{e_1} \oplus \ldots \oplus P_{e_m})
$$
中的像确定。这样便得到单射
$$
T\Deformationcategory_P \longrightarrow W
$$
显然 $W$ 有限维，故引理成立。
\end{proof}

\noindent
在例 \ref{examples-defos-example-graded-algebras} 中，若 $P$ 是有限生成分次
$k$-代数，则 $\Deformationcategory_P$ 可由 $\mathcal{C}_\Lambda$ 上
函子范畴中的光滑 pro-可表示群胚给出表示，因而具有一个（极小）通用
形式对象。这来自引理 \ref{examples-defos-lemma-graded-algebras-RS}、
\ref{examples-defos-lemma-graded-algebras-TI} 以及第 \ref{examples-defos-section-general} 节的一般讨论。

\begin{lemma}
\label{examples-defos-lemma-graded-algebras-hull}
在例 \ref{examples-defos-example-graded-algebras} 中，假设 $P$ 是有限生成分次
$k$-代数，并假设 $\Lambda$ 是剩余域为 $k$ 的完备局部环（经典情形）。
则函子
$$
F : \mathcal{C}_\Lambda \longrightarrow \textit{Sets},\quad
A \longmapsto \Ob(\Deformationcategory_P(A))/\cong
$$
（它取对象的同构类）具有包络。
\end{lemma}

\begin{proof}
这立即来自引理 \ref{examples-defos-lemma-graded-algebras-RS}、
\ref{examples-defos-lemma-graded-algebras-TI}、《形式形变理论》引理
\ref{formal-defos-lemma-RS-implies-S1-S2} 与注
\ref{formal-defos-remark-compose-minimal-into-iso-classes}。
\end{proof}







\section{环}
\label{examples-defos-section-rings}

\noindent
环的形变理论与仿射概形的形变理论相同。对环和概形而言，我们所说的
形变都是指{\it 平坦}形变。

\begin{example}[环]
\label{examples-defos-example-rings}
如下定义范畴 $\mathcal{F}$：
\begin{enumerate}
\item 对象是二元组 $(A,P)$，其中 $A$ 是 $\mathcal{C}_\Lambda$ 的对象，
$P$ 是平坦 $A$-代数；
\item 态射 $(f,g):(B,Q)\to(A,P)$ 由 $\mathcal{C}_\Lambda$ 中的态射
$f:B\to A$ 与映射 $g:Q\to P$ 组成，其中该映射是 $f$-线性的，
并诱导同构 $Q\otimes_{B,f}A\cong P$。
\end{enumerate}
函子 $p:\mathcal{F}\to\mathcal{C}_\Lambda$ 把 $(A,P)$ 送到 $A$，
把 $(f,g)$ 送到 $f$。显然，$p$ 以群胚为纤维余纤维化。
给定 $k$-代数 $P$，令 $x_0=(k,P)$ 为 $\mathcal{F}(k)$ 的相应对象。记
$$
\Deformationcategory_P = \mathcal{F}_{x_0}
$$
\end{example}

\begin{lemma}
\label{examples-defos-lemma-rings-RS}
例 \ref{examples-defos-example-rings} 满足 Rim--Schlessinger 条件 (RS)。特别地，
对任意 $k$-代数 $P$，$\Deformationcategory_P$ 是形变范畴。
\end{lemma}

\begin{proof}
设 $A_1\to A$ 与 $A_2\to A$ 是 $\mathcal{C}_\Lambda$ 的态射，
并假设 $A_2\to A$ 满射。由《形式形变理论》引理
\ref{formal-defos-lemma-RS-2-categorical}，只须证明函子
$\mathcal{F}(A_1 \times_A A_2) \to
\mathcal{F}(A_1) \times_{\mathcal{F}(A)} \mathcal{F}(A_2)$
是范畴等价。这是《更多代数》引理
\ref{more-algebra-lemma-properties-algebras-over-fibre-product}
的特殊情形。
\end{proof}

\begin{lemma}
\label{examples-defos-lemma-rings-TI}
在例 \ref{examples-defos-example-rings} 中，设 $P$ 为 $k$-代数。则
$$
T\Deformationcategory_P = \text{Ext}^1_P(\NL_{P/k}, P)
\quad\text{且}\quad
\text{Inf}(\Deformationcategory_P) = \text{Der}_k(P, P)
$$
\end{lemma}

\begin{proof}
回忆，$\text{Inf}(\Deformationcategory_P)$ 是 $P$ 在 $k[\epsilon]$
上的平凡形变 $P[\epsilon]=P\otimes_k k[\epsilon]$ 中模 $\epsilon$
等于恒等的自同构之集。由《形变理论》引理
\ref{defos-lemma-huge-diagram}，它等于 $\Hom_P(\Omega_{P/k},P)$；
再由《代数》引理 \ref{algebra-lemma-universal-omega}，后者等于
$\text{Der}_k(P,P)$。

\medskip\noindent
回忆，$T\Deformationcategory_P$ 是 $P$ 在 $k[\epsilon]$ 上的平坦形变
$Q$ 的同构类之集；更确切地说，它是
$\Deformationcategory_P(k[\epsilon])$ 的同构类之集。再回忆，满足
$Q/\epsilon Q=P$ 的 $k[\epsilon]$-代数 $Q$ 在 $k[\epsilon]$ 上平坦，
当且仅当
$$
0 \to P \xrightarrow{\epsilon} Q \to P \to 0
$$
正合。这在《更多态射》引理 \ref{more-morphisms-lemma-deform} 中得到证明；
更一般的结果见《形变理论》引理 \ref{defos-lemma-deform-module}。
因此可应用《形变理论》引理 \ref{defos-lemma-choices}，得到这类形变的
同构类之集等于 $\text{Ext}^1_P(\NL_{P/k},P)$。
\end{proof}

\begin{lemma}
\label{examples-defos-lemma-smooth}
在例 \ref{examples-defos-example-rings} 中，设 $P$ 为光滑 $k$-代数。则
$T\Deformationcategory_P=(0)$。
\end{lemma}

\begin{proof}
由引理 \ref{examples-defos-lemma-rings-TI}，须证明
$\text{Ext}^1_P(\NL_{P/k},P)=(0)$。由于 $k\to P$ 光滑，
$\NL_{P/k}$ 拟同构于由置于次数 $0$ 的有限投射 $P$-模组成的复形。
\end{proof}

\begin{lemma}
\label{examples-defos-lemma-finite-type-rings-TI}
在引理 \ref{examples-defos-lemma-rings-TI} 中，若 $P$ 是有限型 $k$-代数，则
\begin{enumerate}
\item $\text{Inf}(\Deformationcategory_P)$ 有限维，当且仅当
$\dim(P)=0$；
\item 若 $\Spec(P)\to\Spec(k)$ 除有限多个点外均光滑，则
$T\Deformationcategory_P$ 有限维。
\end{enumerate}
\end{lemma}

\begin{proof}
证明 (1)。把 $\text{Der}_k(P,P)$ 看作 $P$-模。若它在 $k$ 上有限维，
则作为 $P$-模具有有限长度，因而其支集包含于 $\Spec(P)$ 的有限多个闭点
（《代数》引理 \ref{algebra-lemma-simple-pieces}）。由于
$\text{Der}_k(P,P)=\Hom_P(\Omega_{P/k},P)$，可见
$\text{Der}_k(P, P)_\mathfrak p = \text{Der}_k(P_\mathfrak p, P_\mathfrak p)$
对任意素理想 $\mathfrak p\subset P$ 都成立（这里使用《代数》引理
\ref{algebra-lemma-differentials-localize},
\ref{algebra-lemma-differentials-finitely-presented} 以及
\ref{algebra-lemma-hom-from-finitely-presented}）。设 $\mathfrak p$
为 $P$ 的极小素理想，对应于维数 $d>0$ 的不可约分支。那么
$P_\mathfrak p$ 是本质有限型于 $k$ 的 Artin 局部环，具有剩余域；并且，
例如由《代数》引理 \ref{algebra-lemma-characterize-smooth-over-field}，
$\Omega_{P_\mathfrak p/k}$ 非零。
% P11-E0047：冻结源句子在 “with residue field” 之后缺少应有的说明。
Artin 局部环上的任意非零有限模，都有同构于剩余域的子模与商模。因此
$\text{Der}_k(P_\mathfrak p, P_\mathfrak p) =
\Hom_{P_\mathfrak p}(\Omega_{P_\mathfrak p/k}, P_\mathfrak p)$
也非零。综合以上各点即得 (1)。

\medskip\noindent
证明 (2)。对 $P$ 的素理想 $\mathfrak p$，我们使用
$\NL_{P_\mathfrak p/k} = (\NL_{P/k})_\mathfrak p$
（《代数》引理 \ref{algebra-lemma-localize-NL}）以及
$\text{Ext}_P^1(\NL_{P/k}, P)_\mathfrak p =
\text{Ext}_{P_\mathfrak p}^1(\NL_{P_\mathfrak p/k}, P_\mathfrak p)$
（《更多代数》引理
\ref{more-algebra-lemma-pseudo-coherence-and-base-change-ext}）。
给定素理想 $\mathfrak p\subset P$，环同态 $k\to P$ 在
$\mathfrak p$ 处光滑，当且仅当 $(\NL_{P/k})_\mathfrak p$ 拟同构于
置于次数 $0$ 的有限投射模（这直接来自光滑环同态的定义，也可由更强的
《代数》引理 \ref{algebra-lemma-smooth-at-point} 得到）。

\medskip\noindent
假设 $P$ 在 $k$ 上除有限多个素理想外均光滑。由《代数》引理
\ref{algebra-lemma-finite-type-algebra-finite-nr-primes} 以及这些“坏”素理想
构成 $\Spec(P)$ 的闭子集这一事实，它们都是极大理想
$\mathfrak m_1,\ldots,\mathfrak m_n\subset P$。对
$\mathfrak p\not\in\{\mathfrak m_1,\ldots,\mathfrak m_n\}$，上述结果给出
$\text{Ext}^1_P(\NL_{P/k},P)_\mathfrak p=0$。因此
$\text{Ext}^1_P(\NL_{P/k},P)$ 是有限 $P$-模，其支集包含于
$\{\mathfrak m_1,\ldots,\mathfrak m_r\}$。
% P11-E0048：冻结源先编号至 m_n，此处却写成 m_r。
例如由《代数》命题
\ref{algebra-proposition-minimal-primes-associated-primes}，
$\text{Ext}^1_P(\NL_{P/k},P)$ 在 $k$ 上的维数是若干
$\dim_k\kappa(\mathfrak m_i)$ 的有限整数线性组合；由 Hilbert 零点定理
（《代数》定理 \ref{algebra-theorem-nullstellensatz}），它是有限的。
\end{proof}

\noindent
在例 \ref{examples-defos-example-rings} 中，设 $P$ 为有限型 $k$-代数。
$\Deformationcategory_P$ 可由 $\mathcal{C}_\Lambda$ 上函子范畴中的
光滑 pro-可表示群胚给出表示，当且仅当 $\dim(P)=0$。此外，若
$\Spec(P)\to\Spec(k)$ 只有有限多个奇点，则
$\Deformationcategory_P$ 具有通用形式对象。这来自引理
\ref{examples-defos-lemma-rings-RS}、\ref{examples-defos-lemma-finite-type-rings-TI} 以及第
\ref{examples-defos-section-general} 节的一般讨论。

\begin{lemma}
\label{examples-defos-lemma-rings-hull}
在例 \ref{examples-defos-example-rings} 中，假设 $P$ 是有限型 $k$-代数，且
$\Spec(P)\to\Spec(k)$ 除有限多个点外均光滑。再假设 $\Lambda$
是剩余域为 $k$ 的完备局部环（经典情形）。则函子
$$
F : \mathcal{C}_\Lambda \longrightarrow \textit{Sets},\quad
A \longmapsto \Ob(\Deformationcategory_P(A))/\cong
$$
（它取对象的同构类）具有包络。
\end{lemma}

\begin{proof}
这立即来自引理 \ref{examples-defos-lemma-rings-RS}、
\ref{examples-defos-lemma-finite-type-rings-TI}、《形式形变理论》引理
\ref{formal-defos-lemma-RS-implies-S1-S2} 与注
\ref{formal-defos-remark-compose-minimal-into-iso-classes}。
\end{proof}

\begin{lemma}
\label{examples-defos-lemma-localization}
在例 \ref{examples-defos-example-rings} 中，设 $P$ 为 $k$-代数，并设 $S\subset P$
为乘法子集。存在形变范畴之间的自然函子
$$
\Deformationcategory_P \longrightarrow \Deformationcategory_{S^{-1}P}
$$
\end{lemma}

\begin{proof}
给定 $P$ 的形变，可以将其局部化，从而得到该局部化的形变；这是显然的，
我们建议读者跳过此证明。更确切地说，设
$(A,Q)\to(k,P)$ 是 $\mathcal{F}$ 中的态射，即
$\Deformationcategory_P$ 的对象。令 $S_Q\subset Q$ 为 $S$ 的原像。
因此，$(A,S_Q^{-1}Q)\to(k,S^{-1}P)$ 是
$\Deformationcategory_{S^{-1}P}$ 的所需对象。
% P11-E0049：冻结源在此有悬空的 “Then” 紧接 “Hence”。
\end{proof}

\begin{lemma}
\label{examples-defos-lemma-henselization}
在例 \ref{examples-defos-example-rings} 中，设 $P$ 为 $k$-代数，$J\subset P$ 为理想。
以 $(P^h,J^h)$ 表示二元组 $(P,J)$ 的 Hensel 化。存在形变范畴之间的
自然函子
$$
\Deformationcategory_P \longrightarrow \Deformationcategory_{P^h}
$$
\end{lemma}

\begin{proof}
给定 $P$ 的形变，可以取其 Hensel 化，从而得到 Hensel 化的形变；
这是显然的，我们建议读者跳过此证明。更确切地说，设
$(A,Q)\to(k,P)$ 是 $\mathcal{F}$ 中的态射，即
$\Deformationcategory_P$ 的对象。以 $J_Q\subset Q$ 表示 $J$ 在 $Q$
中的原像，并设 $(Q^h,J_Q^h)$ 为二元组 $(Q,J_Q)$ 的 Hensel 化。
回忆，$Q\to Q^h$ 平坦（《更多代数》引理
\ref{more-algebra-lemma-henselization-flat}），故 $Q^h$ 在 $A$ 上平坦。
由《更多代数》引理 \ref{more-algebra-lemma-henselization-integral}，
态射 $Q^h\to P^h$ 诱导同构
$Q^h\otimes_Ak=Q^h\otimes_QP=P^h$。因此，
$(A,Q^h)\to(k,P^h)$ 是 $\Deformationcategory_{P^h}$ 的所需对象。
\end{proof}

\begin{lemma}
\label{examples-defos-lemma-strict-henselization}
在例 \ref{examples-defos-example-rings} 中，设 $P$ 为 $k$-代数。假设 $P$ 是局部环，
并设 $P^{sh}$ 为 $P$ 的一个严格 Hensel 化。存在形变范畴之间的自然函子
$$
\Deformationcategory_P \longrightarrow \Deformationcategory_{P^{sh}}
$$
\end{lemma}

\begin{proof}
给定 $P$ 的形变，可以取其严格 Hensel 化，从而得到严格 Hensel 化的
形变；这是显然的，我们建议读者跳过此证明。更确切地说，设
$(A,Q)\to(k,P)$ 是 $\mathcal{F}$ 中的态射，即
$\Deformationcategory_P$ 的对象。由于满射 $Q\to P$ 的核幂零，
$Q$ 是与 $P$ 具有相同剩余域的局部环。设 $Q^{sh}$ 为 $Q$ 的严格
Hensel 化。回忆，$Q\to Q^{sh}$ 平坦（《更多代数》引理
\ref{more-algebra-lemma-dumb-properties-henselization}），故 $Q^{sh}$
在 $A$ 上平坦。由《代数》引理
\ref{algebra-lemma-quotient-strict-henselization}，态射
$Q^{sh}\to P^{sh}$ 诱导同构
$Q^{sh}\otimes_Ak=Q^{sh}\otimes_QP=P^{sh}$。因此，
$(A,Q^{sh})\to(k,P^{sh})$ 是
$\Deformationcategory_{P^{sh}}$ 的所需对象。
\end{proof}

\begin{lemma}
\label{examples-defos-lemma-completion}
在例 \ref{examples-defos-example-rings} 中，设 $P$ 为 $k$-代数。假设 $P$ 是 Noether 环，
并设 $J\subset P$ 为理想，以 $P^\wedge$ 表示 $J$-进完备化。
存在形变范畴之间的自然函子
$$
\Deformationcategory_P \longrightarrow \Deformationcategory_{P^\wedge}
$$
\end{lemma}

\begin{proof}
给定 $P$ 的形变，可以将其完备化，从而得到完备化的形变；这是显然的，
我们建议读者跳过此证明。更确切地说，设
$(A,Q)\to(k,P)$ 是 $\mathcal{F}$ 中的态射，即
$\Deformationcategory_P$ 的对象。注意 $Q$ 是诺特环：满环同态
$Q\to P$ 的核幂零且有限生成，而 $P$ 是诺特环；应用《代数》引理
\ref{algebra-lemma-completion-Noetherian}。以 $J_Q\subset Q$ 表示 $J$
在 $Q$ 中的原像，并设 $Q^\wedge$ 为 $Q$ 的 $J_Q$-进完备化。
回忆，$Q\to Q^\wedge$ 平坦（《代数》引理
\ref{algebra-lemma-completion-flat}），故 $Q^\wedge$ 在 $A$ 上平坦。
例如由《代数》引理 \ref{algebra-lemma-completion-tensor}，诱导态射
$Q^\wedge\to P^\wedge$ 诱导同构
$Q^\wedge\otimes_Ak=Q^\wedge\otimes_QP=P^\wedge$。因此，
$(A,Q^\wedge)\to(k,P^\wedge)$ 是
$\Deformationcategory_{P^\wedge}$ 的所需对象。
\end{proof}

\begin{lemma}
\label{examples-defos-lemma-power-series-rings-TI}
在引理 \ref{examples-defos-lemma-rings-TI} 中，若
$P=k[[x_1,\ldots,x_n]]/(f)$，其中非零的
$f\in(x_1,\ldots,x_n)^2$，则
\begin{enumerate}
\item $\text{Inf}(\Deformationcategory_P)$ 有限维，当且仅当 $n=1$；
\item 若
$$
\sqrt{(f, \partial f/\partial x_1, \ldots,  \partial f/\partial x_n)} =
(x_1, \ldots, x_n)
$$
，则 $T\Deformationcategory_P$ 有限维。
\end{enumerate}
\end{lemma}

\begin{proof}
证明 (1)。考虑 $k[[x_1,\ldots,x_n]]$ 在 $k$ 上的导子
$\partial/\partial x_i$。记 $f_i=\partial f/\partial x_i$。导子
$$
\theta = \sum h_i \partial/\partial x_i
$$
是 $k[[x_1,\ldots,x_n]]$ 的导子；它在
$P=k[[x_1,\ldots,x_n]]/(f)$ 上诱导导子，当且仅当
$\sum h_if_i\in(f)$。此外，所诱导的 $P$ 上导子为零，当且仅当对
$i=1,\ldots,n$ 均有 $h_i\in(f)$。因此得到
$$
\Ker((f_1, \ldots, f_n) : P^{\oplus n} \longrightarrow P) \subset
\text{Der}_k(P, P)
$$
左端仅在 $n=1$ 时才是有限维 $k$-向量空间；证明从略。右端在 $n=1$
时有限维这一点也留给读者验证。由此得到 (1)。

\medskip\noindent
证明 (2)。令 $Q$ 为 $P$ 在 $k[\epsilon]$ 上的平坦形变，如引理
\ref{examples-defos-lemma-rings-TI} 的证明中那样。选取 $x_i$ 在 $P$ 中之像的提升
$q_i\in Q$。则 $Q$ 是完备局部环，其极大理想由
$q_1,\ldots,q_n$ 与 $\epsilon$ 生成（略去一个简短论证）。因此得到满射
$$
k[\epsilon][[x_1, \ldots, x_n]] \longrightarrow Q,\quad
x_i \longmapsto q_i
$$
选取映到 $Q$ 中零的元素
$f+\epsilon g\in k[\epsilon][[x_1,\ldots,x_n]]$。注意，$g$ 模 $(f)$
良定义。由于 $Q$ 在 $k[\epsilon]$ 上平坦，得到
$$
Q = k[\epsilon][[x_1, \ldots, x_n]]/(f + \epsilon g)
$$
最后，改变 $q_i$ 的选择，相当于对某些
$h_i\in k[[x_1,\ldots,x_n]]$，把坐标 $x_i$ 改为
$x_i+\epsilon h_i$。于是 $f+\epsilon g$ 变为
$f+\epsilon(g+\sum h_if_i)$，其中 $f_i=\partial f/\partial x_i$。
% P11-E0051：冻结源此句写作 “if we changing”。
因此，形变 $Q$ 的同构类由下列空间中的一个元素确定：
$$
k[[x_1, \ldots, x_n]]/
(f, \partial f/\partial x_1, \ldots,  \partial f/\partial x_n)
$$
它在 $k$ 上有限维，当且仅当其支集是
$k[[x_1,\ldots,x_n]]$ 的闭点；这又当且仅当
$\sqrt{(f, \partial f/\partial x_1, \ldots,  \partial f/\partial x_n)} =
(x_1, \ldots, x_n)$.
\end{proof}






\section{概形}
\label{examples-defos-section-schemes}

\noindent
概形的形变理论。

\begin{example}[概形]
\label{examples-defos-example-schemes}
如下定义范畴 $\mathcal{F}$：
\begin{enumerate}
\item 对象是二元组 $(A,X)$，其中 $A$ 是 $\mathcal{C}_\Lambda$ 的对象，
$X$ 是在 $A$ 上平坦的概形；
\item 态射 $(f,g):(B,Y)\to(A,X)$ 由 $\mathcal{C}_\Lambda$ 中的态射
$f:B\to A$ 与态射 $g:X\to Y$ 组成，使概形的交换图
$$
\xymatrix{
X \ar[r]_g \ar[d] & Y \ar[d] \\
\Spec(A) \ar[r]^f & \Spec(B)
}
$$
为笛卡儿图。
\end{enumerate}
函子 $p:\mathcal{F}\to\mathcal{C}_\Lambda$ 把 $(A,X)$ 送到 $A$，
把 $(f,g)$ 送到 $f$。显然，$p$ 以群胚为纤维余纤维化。
给定 $k$ 上概形 $X$，令 $x_0=(k,X)$ 为 $\mathcal{F}(k)$ 的相应对象。记
$$
\Deformationcategory_X = \mathcal{F}_{x_0}
$$
\end{example}

\begin{lemma}
\label{examples-defos-lemma-schemes-RS}
例 \ref{examples-defos-example-schemes} 满足 Rim--Schlessinger 条件 (RS)。特别地，
对任意 $k$ 上概形 $X$，$\Deformationcategory_X$ 是形变范畴。
\end{lemma}

\begin{proof}
设 $A_1\to A$ 与 $A_2\to A$ 是 $\mathcal{C}_\Lambda$ 的态射，
并假设 $A_2\to A$ 满射。由《形式形变理论》引理
\ref{formal-defos-lemma-RS-2-categorical}，只须证明函子
$\mathcal{F}(A_1 \times_A A_2) \to
\mathcal{F}(A_1) \times_{\mathcal{F}(A)} \mathcal{F}(A_2)$
是范畴等价。注意
$$
\xymatrix{
\Spec(A) \ar[r] \ar[d] & \Spec(A_2) \ar[d] \\
\Spec(A_1) \ar[r] &
\Spec(A_1 \times_A A_2)
}
$$
是《更多态射》引理
\ref{more-morphisms-lemma-pushout-along-thickening} 中的推出图。
因此，本引理是《更多态射》引理
\ref{more-morphisms-lemma-equivalence-categories-schemes-over-pushout-flat}
的特殊情形。
\end{proof}

\begin{lemma}
\label{examples-defos-lemma-schemes-TI}
在例 \ref{examples-defos-example-schemes} 中，设 $X$ 为 $k$ 上概形。则
$$
\text{Inf}(\Deformationcategory_X) =
\text{Ext}^0_{\mathcal{O}_X}(\NL_{X/k}, \mathcal{O}_X) =
\Hom_{\mathcal{O}_X}(\Omega_{X/k}, \mathcal{O}_X) =
\text{Der}_k(\mathcal{O}_X, \mathcal{O}_X)
$$
且
$$
T\Deformationcategory_X =
\text{Ext}^1_{\mathcal{O}_X}(\NL_{X/k}, \mathcal{O}_X)
$$
\end{lemma}

\begin{proof}
回忆，$\text{Inf}(\Deformationcategory_X)$ 是 $X$ 在 $k[\epsilon]$
上的平凡形变
$X'=X\times_{\Spec(k)}\Spec(k[\epsilon])$ 中模 $\epsilon$ 等于恒等的
自同构之集。由《形变理论》引理 \ref{defos-lemma-deform}，它等于
$\text{Ext}^0_{\mathcal{O}_X}(\NL_{X/k},\mathcal{O}_X)$。等式
$\text{Ext}^0_{\mathcal{O}_X}(\NL_{X/k},\mathcal{O}_X)=
\Hom_{\mathcal{O}_X}(\Omega_{X/k},\mathcal{O}_X)$ 来自《更多态射》引理
\ref{more-morphisms-lemma-netherlander-quasi-coherent}。等式
$\Hom_{\mathcal{O}_X}(\Omega_{X/k}, \mathcal{O}_X) =
\text{Der}_k(\mathcal{O}_X, \mathcal{O}_X)$
来自《态射》引理 \ref{morphisms-lemma-universal-derivation-universal}。

\medskip\noindent
回忆，$T_{x_0}\Deformationcategory_X$ 是 $X$ 在 $k[\epsilon]$ 上的
平坦形变 $X'$ 的同构类之集；更确切地说，它是
$\Deformationcategory_X(k[\epsilon])$ 的同构类之集。因此，引理的
第二个断言来自《形变理论》引理 \ref{defos-lemma-deform}。
\end{proof}

\begin{lemma}
\label{examples-defos-lemma-proper-schemes-TI}
在引理 \ref{examples-defos-lemma-schemes-TI} 中，若 $X$ 在 $k$ 上固有，则
$\text{Inf}(\Deformationcategory_X)$ 与 $T\Deformationcategory_X$
均为有限维。
\end{lemma}

\begin{proof}
由该引理，须证明
$\Ext^1_{\mathcal{O}_X}(\NL_{X/k},\mathcal{O}_X)$ 与
$\Ext^0_{\mathcal{O}_X}(\NL_{X/k},\mathcal{O}_X)$ 有限维。
由《更多态射》引理 \ref{more-morphisms-lemma-netherlander-fp}
以及 $X$ 为诺特概形这一事实，$\NL_{X/k}$ 的上同调层均凝聚，
且除次数 $0$ 与 $-1$ 外均为零。由《概形的导出范畴》引理
\ref{perfect-lemma-ext-finite}，所示 $\Ext$ 群都是有限维
$k$-向量空间，证明完毕。
\end{proof}

\noindent
在例 \ref{examples-defos-example-schemes} 中，若 $X$ 是 $k$ 上固有概形，则
$\Deformationcategory_X$ 可由 $\mathcal{C}_\Lambda$ 上函子范畴中的
光滑 pro-可表示群胚给出表示，因而具有一个（极小）通用形式对象。
这来自引理 \ref{examples-defos-lemma-schemes-RS}、\ref{examples-defos-lemma-proper-schemes-TI}
以及第 \ref{examples-defos-section-general} 节的一般讨论。

\begin{lemma}
\label{examples-defos-lemma-schemes-hull}
在例 \ref{examples-defos-example-schemes} 中，假设 $X$ 是固有 $k$-概形，并假设
$\Lambda$ 是剩余域为 $k$ 的完备局部环（经典情形）。则函子
$$
F : \mathcal{C}_\Lambda \longrightarrow \textit{Sets},\quad
A \longmapsto \Ob(\Deformationcategory_X(A))/\cong
$$
（它取对象的同构类）具有包络。若
$\text{Der}_k(\mathcal{O}_X,\mathcal{O}_X)=0$，则 $F$ 是
pro-可表示的。
\end{lemma}

\begin{proof}
包络的存在性立即来自引理 \ref{examples-defos-lemma-schemes-RS}、
\ref{examples-defos-lemma-proper-schemes-TI}、《形式形变理论》引理
\ref{formal-defos-lemma-RS-implies-S1-S2} 与注
\ref{formal-defos-remark-compose-minimal-into-iso-classes}。

\medskip\noindent
假设 $\text{Der}_k(\mathcal{O}_X,\mathcal{O}_X)=0$。由《形式形变理论》
引理 \ref{formal-defos-lemma-infdef-trivial}，
$\Deformationcategory_X$ 与 $F$ 等价。因此，$F$ 是具有有限维切空间的
形变函子（因为 $\Deformationcategory_X$ 是形变范畴），可以应用
《形式形变理论》定理
\ref{formal-defos-theorem-Schlessinger-prorepresentability}。
\end{proof}

\begin{lemma}
\label{examples-defos-lemma-open}
在例 \ref{examples-defos-example-schemes} 中，设 $X$ 为 $k$ 上概形，$U\subset X$
为开子概形。存在形变范畴之间的自然函子
$$
\Deformationcategory_X \longrightarrow \Deformationcategory_U
$$
\end{lemma}

\begin{proof}
给定 $X$ 的形变，可以取其中相应的开子概形，从而得到 $U$ 的形变。
细节从略。
\end{proof}

\begin{lemma}
\label{examples-defos-lemma-affine}
在例 \ref{examples-defos-example-schemes} 中，设 $X=\Spec(P)$ 为仿射 $k$-概形。
令 $\Deformationcategory_P$ 如例 \ref{examples-defos-example-rings}。存在形变范畴的
自然等价
$$
\Deformationcategory_X \longrightarrow \Deformationcategory_P
$$
\end{lemma}

\begin{proof}
该函子把 $(A,Y)$ 送到 $\Gamma(Y,\mathcal{O}_Y)$。这是可行的，因为由
《更多态射》引理 \ref{more-morphisms-lemma-thickening-affine-scheme}，
$X$ 的任意形变都是仿射的。
\end{proof}

\begin{lemma}
\label{examples-defos-lemma-local-ring}
在例 \ref{examples-defos-example-schemes} 中，设 $X$ 为 $k$ 上概形，并设 $p\in X$
为一点。令 $\Deformationcategory_{\mathcal{O}_{X,p}}$ 如例
\ref{examples-defos-example-rings}。存在形变范畴之间的自然函子
$$
\Deformationcategory_X
\longrightarrow
\Deformationcategory_{\mathcal{O}_{X, p}}
$$
\end{lemma}

\begin{proof}
选取包含 $p$ 的仿射开子概形 $U=\Spec(P)\subset X$。则
$\mathcal{O}_{X,p}$ 是 $P$ 的一个局部化。复合引理
\ref{examples-defos-lemma-open}、\ref{examples-defos-lemma-affine} 与 \ref{examples-defos-lemma-localization}
中的函子即可。
\end{proof}

\begin{situation}
\label{examples-defos-situation-glueing}
设 $\Lambda\to k$ 如第 \ref{examples-defos-section-general} 节。设 $X$ 为 $k$ 上概形，
具有仿射开覆盖 $X=U_1\cup U_2$，且
$U_{12}=U_1\cap U_2$ 也仿射。写成
$U_1=\Spec(P_1)$、$U_2=\Spec(P_2)$ 与 $U_{12}=\Spec(P_{12})$。
令 $\Deformationcategory_X$、$\Deformationcategory_{U_1}$、
$\Deformationcategory_{U_2}$ 与 $\Deformationcategory_{U_{12}}$
如例 \ref{examples-defos-example-schemes}，并令
$\Deformationcategory_{P_1}$、$\Deformationcategory_{P_2}$ 与
$\Deformationcategory_{P_{12}}$ 如例 \ref{examples-defos-example-rings}。
\end{situation}

\begin{lemma}
\label{examples-defos-lemma-glueing}
在情形 \ref{examples-defos-situation-glueing} 中，存在形变范畴的等价
$$
\Deformationcategory_X =
\Deformationcategory_{P_1}
\times_{\Deformationcategory_{P_{12}}}
\Deformationcategory_{P_2}
$$
，参见例 \ref{examples-defos-example-schemes} 与 \ref{examples-defos-example-rings}。
\end{lemma}

\begin{proof}
只须证明引理 \ref{examples-defos-lemma-open} 的函子定义等价
$$
\Deformationcategory_X \longrightarrow
\Deformationcategory_{U_1}
\times_{\Deformationcategory_{U_{12}}}
\Deformationcategory_{U_2}
$$
，因为这样便可应用引理 \ref{examples-defos-lemma-affine} 转化为环的情形。
为此构造一个拟逆。以
$F_i:\Deformationcategory_{U_i}\to\Deformationcategory_{U_{12}}$
表示引理 \ref{examples-defos-lemma-open} 的函子。右端的对象由
$\mathcal{C}_\Lambda$ 中的 $A$、对象
$(A,V_1)\to(k,U_1)$ 与 $(A,V_2)\to(k,U_2)$，以及态射
$$
g : F_1(A, V_1) \to F_2(A, V_2)
$$
给出。这里 $F_i(A,V_i)=(A,V_{i,3-i})$，其中
$V_{i,3-i}\subset V_i$ 是基变换到 $k$ 后成为
$U_{12}\subset U_i$ 的开子概形。态射 $g$ 定义 $A$ 上概形的同构
$V_{1,2}\to V_{2,1}$，且它与 $k$ 上的
$\text{id}:U_{12}\to U_{12}$ 相容。因此
$(\{1,2\},V_i,V_{i,3-i},g,g^{-1})$ 是《概形》第
\ref{schemes-section-glueing-schemes} 节中的粘合数据。设 $Y$ 为所得粘合，
参见《概形》引理 \ref{schemes-lemma-glue}。则 $Y$ 是 $A$ 上概形，
上述相容性给出典范同构
$Y\times_{\Spec(A)}\Spec(k)=X$。因此
$(A,Y)\to(k,X)$ 是 $\Deformationcategory_X$ 的对象。
关于此构造是函子且为给定函子的拟逆这一验证从略。
\end{proof}






\section{概形态射}
\label{examples-defos-section-schemes-morphisms}

\noindent
本节讨论概形态射的形变理论。
当然，这不过是概形图表形变的一个例子。

\begin{example}[概形态射]
\label{examples-defos-example-schemes-morphisms}
定义范畴 $\mathcal{F}$ 如下：
\begin{enumerate}
\item 对象是一个偶对 $(A, X \to Y)$，其中 $A$ 是
$\mathcal{C}_\Lambda$ 的对象，而 $X \to Y$ 是 $A$ 上的概形态射，
且 $X$ 与 $Y$ 都在 $A$ 上平坦；
\item 态射 $(f, g, h) : (A', X' \to Y') \to (A, X \to Y)$
由 $\mathcal{C}_\Lambda$ 中的态射 $f : A' \to A$，以及概形态射
$g : X \to X'$ 与 $h : Y \to Y'$ 组成，并使得
$$
\xymatrix{
X \ar[r]_g \ar[d] & X' \ar[d] \\
Y \ar[r]_h \ar[d] & Y' \ar[d] \\
\Spec(A) \ar[r]^f & \Spec(A')
}
$$
是概形的交换图，且其中两个方块都是笛卡儿的。
\end{enumerate}
函子 $p : \mathcal{F} \to \mathcal{C}_\Lambda$ 把 $(A, X \to Y)$
送到 $A$，把 $(f, g, h)$ 送到 $f$。显然，$p$ 以群胚为纤维余纤维化。
给定 $k$ 上的概形态射 $X \to Y$，令 $x_0 = (k, X \to Y)$
为 $\mathcal{F}(k)$ 中相应的对象。置
$$
\Deformationcategory_{X \to Y} = \mathcal{F}_{x_0}
$$
\end{example}

\begin{lemma}
\label{examples-defos-lemma-schemes-morphisms-RS}
例 \ref{examples-defos-example-schemes-morphisms}
满足 Rim--Schlessinger 条件 (RS)。
特别地，对 $k$ 上任意概形态射 $X \to Y$，
$\Deformationcategory_{X \to Y}$ 都是形变范畴。
\end{lemma}

\begin{proof}
设 $A_1 \to A$ 与 $A_2 \to A$ 是 $\mathcal{C}_\Lambda$ 中的态射，
并假设 $A_2 \to A$ 是满射。根据形式形变理论引理
\ref{formal-defos-lemma-RS-2-categorical}
，只需证明函子
$\mathcal{F}(A_1 \times_A A_2) \to
\mathcal{F}(A_1) \times_{\mathcal{F}(A)} \mathcal{F}(A_2)$
是范畴等价。注意到
$$
\xymatrix{
\Spec(A) \ar[r] \ar[d] & \Spec(A_2) \ar[d] \\
\Spec(A_1) \ar[r] &
\Spec(A_1 \times_A A_2)
}
$$
是态射进阶引理
\ref{more-morphisms-lemma-pushout-along-thickening} 中的推出图。
因而本引理由态射进阶引理
\ref{more-morphisms-lemma-equivalence-categories-schemes-over-pushout-flat}
立即得到：该引理把在 $A_1 \times_A A_2$ 上平坦的概形所成范畴描述为
在 $A_1$ 上平坦的概形所成范畴与在 $A_2$ 上平坦的概形所成范畴，
在 $A$ 上平坦的概形所成范畴之上的纤维积。
\end{proof}

\begin{lemma}
\label{examples-defos-lemma-schemes-morphisms-TI}
% P11-E0052：冻结源在本节两处引用 example-schemes，而所述对象属于 example-schemes-morphisms。
在例 \ref{examples-defos-example-schemes} 中，设 $f : X \to Y$ 是 $k$ 上的概形态射。
有如下典范的 $k$-向量空间正合列：
$$
\xymatrix{
0 \ar[r] &
\text{Inf}(\Deformationcategory_{X \to Y}) \ar[r] &
\text{Inf}(\Deformationcategory_X \times \Deformationcategory_Y) \ar[r] &
\text{Der}_k(\mathcal{O}_Y, f_*\mathcal{O}_X) \ar[lld] \\
& T\Deformationcategory_{X \to Y} \ar[r] &
T(\Deformationcategory_X \times \Deformationcategory_Y) \ar[r] &
\text{Ext}^1_{\mathcal{O}_X}(Lf^*\NL_{Y/k}, \mathcal{O}_X)
}
$$
\end{lemma}

\begin{proof}
形变范畴间的显然映射
$\Deformationcategory_{X \to Y} \to
\Deformationcategory_X \times \Deformationcategory_Y$
给出了引理中正合列的两个箭头。回忆，
$\text{Inf}(\Deformationcategory_{X \to Y})$
是平凡形变
$$
f' :  X' = X \times_{\Spec(k)} \Spec(k[\epsilon])
\xrightarrow{f \times \text{id}}
Y' = Y \times_{\Spec(k)} \Spec(k[\epsilon])
$$
在模 $\epsilon$ 意义下等于恒等的自同构所成的集合；这里该平凡形变是
$X \to Y$ 到 $k[\epsilon]$ 上的形变。这显然等同于偶对
$(\alpha, \beta) \in
\text{Inf}(\Deformationcategory_X \times \Deformationcategory_Y)$
，其中二者分别是 $X$ 与 $Y$ 的无穷小自同构，
且与 $f'$ 相容；也就是说 $f' \circ \alpha = \beta \circ f'$。
由形变理论引理 \ref{defos-lemma-huge-diagram-ringed-spaces}，
对任意偶对 $(\alpha, \beta)$，态射 $f' : X' \to Y'$ 与态射
$\beta^{-1} \circ f' \circ \alpha : X' \to Y'$ 之差定义了
$$
\text{Der}_k(\mathcal{O}_Y, f_*\mathcal{O}_X) =
\Hom_{\mathcal{O}_Y}(\Omega_{Y/k}, f_*\mathcal{O}_X)
$$
中的一个元素。这里的等号由态射进阶引理
\ref{more-morphisms-lemma-netherlander-quasi-coherent} 给出。
这定义了最上方最后一个水平箭头，并证明了前两处的正合性。对于映射
$$
\text{Der}_k(\mathcal{O}_Y, f_*\mathcal{O}_X)
\to
T\Deformationcategory_{X \to Y}
$$
，利用形变理论引理 \ref{defos-lemma-huge-diagram-ringed-spaces}，
我们把源中的元素解释为 $\Spec(k[\epsilon])$ 上的态射
$f_\epsilon : X' \to Y'$，并要求其模 $\epsilon$ 等于 $f$。
我们把 $f_\epsilon$ 送到
$T\Deformationcategory_{X \to Y}$ 中
$(f_\epsilon : X' \to Y')$ 的同构类。注意，
$(f_\epsilon : X' \to Y')$ 与平凡形变
$(f' : X' \to Y')$ 同构，当且仅当
% P11-E0053：冻结源下一公式中的 f 按类型应为 f'。
对某个偶对 $(\alpha, \beta)$ 有
$f_\epsilon  = \beta^{-1} \circ f \circ \alpha$；
这说明第三处正合。显然，若某个一阶形变
$(f_\epsilon : X_\epsilon \to Y_\epsilon)$
在 $T(\Deformationcategory_X \times \Deformationcategory_Y)$ 中映为零，
则可选取同构 $X' \to X_\epsilon$ 与 $Y' \to Y_\epsilon$，
从而得到它位于左下方向箭头的像中。因此第四处正合。
最后，给定 $X$、$Y$ 的两个一阶形变 $X_\epsilon$、$Y_\epsilon$，
存在障碍
$$
ob(X_\epsilon, Y_\epsilon) \in
\text{Ext}^1_{\mathcal{O}_X}(Lf^*\NL_{Y/k}, \mathcal{O}_X)
$$
；该障碍消失，当且仅当 $f : X \to Y$ 可提升为
$X_\epsilon \to Y_\epsilon$；参见形变理论引理
\ref{defos-lemma-huge-diagram-ringed-spaces}。证明完毕。
\end{proof}

\begin{lemma}
\label{examples-defos-lemma-proper-schemes-morphisms-TI}
在引理 \ref{examples-defos-lemma-schemes-morphisms-TI} 中，若 $X$ 与 $Y$
都在 $k$ 上固有，则
$\text{Inf}(\Deformationcategory_{X \to Y})$ 与
$T\Deformationcategory_{X \to Y}$ 都是有限维的。
\end{lemma}

\begin{proof}
从略。提示：仿照引理 \ref{examples-defos-lemma-proper-schemes-TI} 论证，
并使用本引理的正合列。
\end{proof}

\noindent
在例 \ref{examples-defos-example-schemes-morphisms} 中，若 $X \to Y$ 是
$k$ 上固有概形之间的态射，则
$\Deformationcategory_{X \to Y}$ 可由 $\mathcal{C}_\Lambda$
上函子中的一个光滑 pro-可表示群胚给出呈示，因而更有一个
（极小）通用形式对象。这由引理 \ref{examples-defos-lemma-schemes-morphisms-RS}、
\ref{examples-defos-lemma-proper-schemes-morphisms-TI} 以及
第 \ref{examples-defos-section-general} 节的一般讨论推出。

\begin{lemma}
\label{examples-defos-lemma-schemes-morphisms-hull}
在例 \ref{examples-defos-example-schemes-morphisms} 中，假设 $X \to Y$
是固有 $k$-概形之间的态射。再假设 $\Lambda$ 是剩余域为 $k$
的完备局部环（经典情形）。则函子
$$
F : \mathcal{C}_\Lambda \longrightarrow \textit{Sets},\quad
A \longmapsto \Ob(\Deformationcategory_{X \to Y}(A))/\cong
$$
把对象映到其同构类；它有包络。若
$\text{Der}_k(\mathcal{O}_X, \mathcal{O}_X) =
\text{Der}_k(\mathcal{O}_Y, \mathcal{O}_Y) = 0$，
则 $F$ 是 pro-可表示的。
\end{lemma}

\begin{proof}
包络的存在性立即由引理 \ref{examples-defos-lemma-schemes-morphisms-RS}、
\ref{examples-defos-lemma-proper-schemes-morphisms-TI}、形式形变理论引理
\ref{formal-defos-lemma-RS-implies-S1-S2} 以及注
\ref{formal-defos-remark-compose-minimal-into-iso-classes} 推出。

\medskip\noindent
假设 $\text{Der}_k(\mathcal{O}_X, \mathcal{O}_X) =
\text{Der}_k(\mathcal{O}_Y, \mathcal{O}_Y) = 0$。
引理 \ref{examples-defos-lemma-schemes-morphisms-TI} 的正合列结合引理
\ref{examples-defos-lemma-schemes-TI} 表明
$\text{Inf}(\Deformationcategory_{X \to Y}) = 0$。
于是由形式形变理论引理 \ref{formal-defos-lemma-infdef-trivial}，
$\Deformationcategory_{X \to Y}$ 与 $F$ 等价。因此 $F$ 是具有
有限维切空间的形变函子（因为 $\Deformationcategory_{X \to Y}$
是形变范畴），我们可以应用形式形变理论定理
\ref{formal-defos-theorem-Schlessinger-prorepresentability}。
\end{proof}

\begin{lemma}
\label{examples-defos-lemma-schemes-morphisms-smooth-to-base}
\begin{reference}
这在 \cite[Section 5.3]{Ravi-Murphys-Law} 与
\cite[Theorem 3.3]{Ran-deformations} 中有所讨论。
\end{reference}
在例 \ref{examples-defos-example-schemes} 中，设 $f : X \to Y$ 是 $k$ 上的概形态射。
若 $f_*\mathcal{O}_X = \mathcal{O}_Y$ 且 $R^1f_*\mathcal{O}_X = 0$，
则形变范畴的态射
$$
\Deformationcategory_{X \to Y} \to \Deformationcategory_X
$$
是等价。
\end{lemma}

\begin{proof}
我们构造引理中遗忘函子的一个拟逆。设 $(A,U)$ 是
$\Deformationcategory_X$ 的对象。给定映射 $X \to U$ 是有限阶加厚，
可据此把 $U$ 与 $X$ 的底拓扑空间等同起来；参见态射进阶第
\ref{more-morphisms-section-thickenings} 节。因此我们可以并确实把
$\mathcal{O}_U$ 看作 $X$ 上的 $A$-代数层；此外，
$U \to \Spec(A)$ 平坦意味着 $\mathcal{O}_U$ 作为 $A$-模层平坦。
特别地，有滤过
$$
0 = \mathfrak m_A^n\mathcal{O}_U \subset
\mathfrak m_A^{n - 1}\mathcal{O}_U \subset \ldots \subset
\mathfrak m_A^2\mathcal{O}_U \subset
\mathfrak m_A\mathcal{O}_U \subset \mathcal{O}_U
$$
；由平坦性，其相继商等于
$\mathcal{O}_X \otimes_k \mathfrak m_A^i/\mathfrak m_A^{i + 1}$
；参见态射进阶引理 \ref{more-morphisms-lemma-deform}，
或更一般的形变理论引理 \ref{defos-lemma-deform-module}。置
$$
\mathcal{O}_V = f_*\mathcal{O}_U
$$
，把它视作 $Y$ 上的 $A$-代数层。由于
$R^1f_*\mathcal{O}_X = 0$，由上述描述可知
$R^1f_*(\mathfrak m_A^i\mathcal{O}_U/\mathfrak m_A^{i + 1}\mathcal{O}_U) = 0$
对所有 $i$ 成立。这意味着序列
$$
0 \to
(f_*\mathcal{O}_X) \otimes_k \mathfrak m_A^i/\mathfrak m_A^{i + 1} \to
f_*(\mathcal{O}_U/\mathfrak m_A^{i + 1}\mathcal{O}_U) \to
f_*(\mathcal{O}_U/\mathfrak m_A^i\mathcal{O}_U) \to 0
$$
对所有 $i$ 都正合。反向应用上述文献（并使用归纳法），可知
$\mathcal{O}_V$ 是平坦的 $A$-代数层，且
$\mathcal{O}_V/\mathfrak m_A\mathcal{O}_V = \mathcal{O}_Y$。
应用态射进阶引理
\ref{more-morphisms-lemma-first-order-thickening}
可知 $(Y,\mathcal{O}_V)$ 是概形，记为 $V$。
等式 $\mathcal{O}_V = f_*\mathcal{O}_U$ 定义了带环空间的态射
$U \to V$，容易看出它是概形态射。结合已经证明的平坦性，
这就完成了证明。
\end{proof}







\section{代数空间}
\label{examples-defos-section-algebraic-spaces}

\noindent
本节讨论代数空间的形变理论。

\begin{example}[代数空间]
\label{examples-defos-example-spaces}
定义范畴 $\mathcal{F}$ 如下：
\begin{enumerate}
\item 对象是偶对 $(A,X)$，其中 $A$ 是 $\mathcal{C}_\Lambda$
的对象，$X$ 是在 $A$ 上平坦的代数空间；
\item 态射 $(f,g):(B,Y)\to(A,X)$ 由
$\mathcal{C}_\Lambda$ 中的态射 $f:B\to A$ 以及
$\Lambda$ 上代数空间的态射 $g:X\to Y$ 组成，并使得
$$
\xymatrix{
X \ar[r]_g \ar[d] & Y \ar[d] \\
\Spec(A) \ar[r]^f & \Spec(B)
}
$$
是代数空间的笛卡儿交换图。
\end{enumerate}
函子 $p:\mathcal{F}\to\mathcal{C}_\Lambda$ 把 $(A,X)$ 送到 $A$，
把 $(f,g)$ 送到 $f$。显然，$p$ 以群胚为纤维余纤维化。
给定 $k$ 上的代数空间 $X$，令 $x_0=(k,X)$ 为
$\mathcal{F}(k)$ 中相应的对象。置
$$
\Deformationcategory_X = \mathcal{F}_{x_0}
$$
\end{example}

\begin{lemma}
\label{examples-defos-lemma-spaces-RS}
例 \ref{examples-defos-example-spaces} 满足 Rim--Schlessinger 条件 (RS)。
特别地，对 $k$ 上任意代数空间 $X$，
$\Deformationcategory_X$ 都是形变范畴。
\end{lemma}

\begin{proof}
设 $A_1\to A$ 与 $A_2\to A$ 是 $\mathcal{C}_\Lambda$ 中的态射，
并假设 $A_2\to A$ 为满射。根据形式形变理论引理
\ref{formal-defos-lemma-RS-2-categorical}
，只需证明函子
$\mathcal{F}(A_1 \times_A A_2) \to
\mathcal{F}(A_1) \times_{\mathcal{F}(A)} \mathcal{F}(A_2)$
是范畴等价。注意到
$$
\xymatrix{
\Spec(A) \ar[r] \ar[d] & \Spec(A_2) \ar[d] \\
\Spec(A_1) \ar[r] &
\Spec(A_1 \times_A A_2)
}
$$
是空间的推出引理
\ref{spaces-pushouts-lemma-pushout-along-thickening} 中的推出图。
因而本引理是空间的推出引理
\ref{spaces-pushouts-lemma-equivalence-categories-spaces-pushout-flat}
的特殊情形。
\end{proof}

\begin{lemma}
\label{examples-defos-lemma-spaces-TI}
在例 \ref{examples-defos-example-spaces} 中，设 $X$ 是 $k$ 上的代数空间。则
$$
\text{Inf}(\Deformationcategory_X) =
\text{Ext}^0_{\mathcal{O}_X}(\NL_{X/k}, \mathcal{O}_X) =
\Hom_{\mathcal{O}_X}(\Omega_{X/k}, \mathcal{O}_X) =
\text{Der}_k(\mathcal{O}_X, \mathcal{O}_X)
$$
以及
$$
T\Deformationcategory_X =
\text{Ext}^1_{\mathcal{O}_X}(\NL_{X/k}, \mathcal{O}_X)
$$
\end{lemma}

\begin{proof}
回忆，$\text{Inf}(\Deformationcategory_X)$ 是 $X$ 到 $k[\epsilon]$
上的平凡形变 $X'=X\times_{\Spec(k)}\Spec(k[\epsilon])$
在模 $\epsilon$ 意义下等于恒等的自同构所成的集合。
由形变理论引理 \ref{defos-lemma-deform-spaces}，它等于
$\text{Ext}^0_{\mathcal{O}_X}(\NL_{X/k},\mathcal{O}_X)$。
等式 $\text{Ext}^0_{\mathcal{O}_X}(\NL_{X/k}, \mathcal{O}_X) =
\Hom_{\mathcal{O}_X}(\Omega_{X/k}, \mathcal{O}_X)$
由空间态射进阶引理
\ref{spaces-more-morphisms-lemma-netherlander-quasi-coherent}
得到。等式
$\Hom_{\mathcal{O}_X}(\Omega_{X/k}, \mathcal{O}_X) =
\text{Der}_k(\mathcal{O}_X, \mathcal{O}_X)$
由空间态射进阶定义
\ref{spaces-more-morphisms-definition-sheaf-differentials}
与位点上的模定义 \ref{sites-modules-definition-module-differentials} 得到。

\medskip\noindent
回忆，$T_{x_0}\Deformationcategory_X$ 是 $X$ 到 $k[\epsilon]$
上的平坦形变 $X'$ 的同构类所成的集合；更准确地说，
它是 $\Deformationcategory_X(k[\epsilon])$ 中同构类所成的集合。
因此引理的第二个断言由形变理论引理
\ref{defos-lemma-deform-spaces} 得到。
\end{proof}

\begin{lemma}
\label{examples-defos-lemma-proper-spaces-TI}
在引理 \ref{examples-defos-lemma-spaces-TI} 中，若 $X$ 在 $k$ 上固有，则
$\text{Inf}(\Deformationcategory_X)$ 与 $T\Deformationcategory_X$
都是有限维的。
\end{lemma}

\begin{proof}
由该引理，我们须证明
$\Ext^1_{\mathcal{O}_X}(\NL_{X/k}, \mathcal{O}_X)$ 与
$\Ext^0_{\mathcal{O}_X}(\NL_{X/k}, \mathcal{O}_X)$ 是有限维的。
由空间态射进阶引理
\ref{spaces-more-morphisms-lemma-netherlander-fp}
以及 $X$ 是 Noether 的这一事实可知，$\NL_{X/k}$ 的凝聚上同调层
除次数 $0$ 与 $-1$ 外均为零。由空间的导出范畴引理
\ref{spaces-perfect-lemma-ext-finite}，上述 $\Ext$ 群都是有限维
$k$-向量空间，证明完毕。
\end{proof}

\noindent
在例 \ref{examples-defos-example-spaces} 中，若 $X$ 是 $k$ 上的固有代数空间，
则 $\Deformationcategory_X$ 可由 $\mathcal{C}_\Lambda$
上函子中的一个光滑 pro-可表示群胚给出呈示，因而更有一个
（极小）通用形式对象。这由引理 \ref{examples-defos-lemma-spaces-RS}、
\ref{examples-defos-lemma-proper-spaces-TI} 以及第 \ref{examples-defos-section-general} 节的一般讨论推出。

\begin{lemma}
\label{examples-defos-lemma-spaces-hull}
在例 \ref{examples-defos-example-spaces} 中，假设 $X$ 是 $k$ 上的固有代数空间。
再假设 $\Lambda$ 是剩余域为 $k$ 的完备局部环（经典情形）。则函子
$$
F : \mathcal{C}_\Lambda \longrightarrow \textit{Sets},\quad
A \longmapsto \Ob(\Deformationcategory_X(A))/\cong
$$
把对象映到其同构类；它有包络。若
$\text{Der}_k(\mathcal{O}_X, \mathcal{O}_X)=0$，
则 $F$ 是 pro-可表示的。
\end{lemma}

\begin{proof}
包络的存在性立即由引理 \ref{examples-defos-lemma-spaces-RS}、
\ref{examples-defos-lemma-proper-spaces-TI}、形式形变理论引理
\ref{formal-defos-lemma-RS-implies-S1-S2} 以及注
\ref{formal-defos-remark-compose-minimal-into-iso-classes} 推出。

\medskip\noindent
假设 $\text{Der}_k(\mathcal{O}_X,\mathcal{O}_X)=0$。
由形式形变理论引理 \ref{formal-defos-lemma-infdef-trivial}，
$\Deformationcategory_X$ 与 $F$ 等价。因此 $F$ 是具有有限维切空间的
形变函子（因为 $\Deformationcategory_X$ 是形变范畴），
我们可以应用形式形变理论定理
\ref{formal-defos-theorem-Schlessinger-prorepresentability}。
\end{proof}





\section{完备化的形变}
\label{examples-defos-section-compare}

\noindent
本节比较由一个代数及其完备化所给出的形变问题。
我们先讨论“可提升性”。

\begin{lemma}
\label{examples-defos-lemma-lift-equivalence-module-derived}
设 $A'\to A$ 是核为幂零理想的环满射，
$A'\to P'$ 是平坦环同态。置 $P=P'\otimes_{A'}A$。
设 $M$ 是在 $A$ 上平坦的 $P$-模。则下列条件等价：
\begin{enumerate}
\item 存在在 $A'$ 上平坦的 $P'$-模 $M'$，使得
$M'\otimes_{P'}P=M$；
\item 存在对象 $K'\in D^-(P')$，使得
$K'\otimes_{P'}^\mathbf{L}P=M$。
\end{enumerate}
\end{lemma}

\begin{proof}
假设 $M'$ 如 (1) 所述。则
% P11-E0054：冻结源下列张量等式链中前、后三个底环的方向按类型不相容。
$$
M = M' \otimes_P P' = M' \otimes_{A'} A =
M' \otimes_A^\mathbf{L} A' = M' \otimes_{P'}^\mathbf{L} P
$$
前两个等式是显然的；第三个等式由 $M'$ 在 $A'$ 上平坦得到；
第四个等式由代数进阶引理
\ref{more-algebra-lemma-base-change-comparison} 得到。因此 (2) 成立。
反之，假设 $K'$ 如 (2) 所述。我们可以并确实假设 $M$ 非零。
令 $t$ 为使 $H^t(K')$ 非零的最大整数（它因 $M$ 非零而存在）。
若 $t>0$，则
$H^t(K')\otimes_{P'}P=H^t(K'\otimes_{P'}^\mathbf{L}P)$ 为零。
由于 $P'\to P$ 的核幂零，由 Nakayama 引理可得 $H^t(K')=0$，
矛盾。因此 $t=0$（$t<0$ 的情形同样荒谬）。于是
$M'=H^0(K')$ 是满足 $M=M'\otimes_{P'}P$ 的 $P'$-模，
而 Tor 谱序列给出单射
$$
\text{Tor}_1^{P'}(M', P) \to H^{-1}(M' \otimes_{P'}^\mathbf{L} P) = 0
$$
。由上述导出基变换文献，
$0 = \text{Tor}_1^{P'}(M', P) = \text{Tor}_1^{A'}(M', A)$.
再由代数引理 \ref{algebra-lemma-what-does-it-mean} 可知，
$M'$ 在 $A'$ 上平坦。
\end{proof}

\begin{lemma}
\label{examples-defos-lemma-lift-equivalence-module}
考虑 Noether 环的交换图
$$
\xymatrix{
A' \ar[d] \ar[r] &
P' \ar[d] \ar[r] &
Q' \ar[d] \\
A \ar[r] &
P \ar[r] &
Q
}
$$
，其中两个方块均为笛卡儿方块，水平箭头均平坦，
竖直箭头均为核幂零的满射。设 $J'\subset P'$ 是满足
$P'/J'=Q'/J'Q'$ 的理想，且 $M$ 是在 $A$ 上平坦的 $P$-模。
假设对每个 $g\in J'$，都存在在 $A'$ 上平坦的
$(P')_g$-模提升 $M_g$。则下列条件等价：
\begin{enumerate}
\item $M$ 有一个在 $A'$ 上平坦的 $P'$-模提升；
\item $M\otimes_PQ$ 有一个在 $A'$ 上平坦的 $Q'$-模提升。
\end{enumerate}
\end{lemma}

\begin{proof}
令 $I=\Ker(A'\to A)$。对满足 $I^n=0$ 的整数 $n>1$ 作归纳，
可约化到 $I$ 是平方零理想的情形；细节从略。
我们把 $M$ 的可提升条件转化为寻找引理
\ref{examples-defos-lemma-lift-equivalence-module-derived} 中那样的
$D^-(P')$ 对象的问题。相应的障碍是元素
$$
\omega(M) \in \text{Ext}^2_P(M, M \otimes_P^\mathbf{L} IP) =
\text{Ext}^2_P(M, M \otimes_P IP)
$$
，它由形变理论引理 \ref{defos-lemma-canonical-class-algebra} 构造。
显示公式中的等式成立，是因为 $M$ 与 $P$ 都在 $A$ 上平坦，故
$M\otimes_P^\mathbf{L}IP=M\otimes_PIP$\footnote{取由自由
$A$-模组成的解消 $F_\bullet\to I$。由于 $A\to P$ 平坦，
$P\otimes_AF_\bullet$ 是 $IP$ 的自由解消。因此
$M\otimes_P^\mathbf{L}IP$ 由
$M\otimes_PP\otimes_AF_\bullet=M\otimes_AF_\bullet$ 表示。
由于 $M$ 在 $A$ 上平坦，该复形只在次数 $0$ 有上同调。}。
同样，提升 $M\otimes_PQ$ 的障碍是元素
$$
\omega(M \otimes_P Q) \in
\text{Ext}^2_Q(M \otimes_P Q, (M \otimes_P Q) \otimes_Q IQ)
$$
；由形变理论引理 \ref{defos-lemma-canonical-class-algebra}
中构造 $\omega(-)$ 的函子性，它是 $\omega(M)$ 的像。
由代数进阶引理 \ref{more-algebra-lemma-base-change-RHom}，有
$$
\text{Ext}^2_Q(M \otimes_P Q, (M \otimes_P Q) \otimes_Q IQ) =
\text{Ext}^2_P(M, M \otimes_P IP) \otimes_P Q
$$
。这里使用了 $P$ 是 Noether 环且 $M$ 有限。
% P11-E0055：冻结源引理的假设并未要求 M 有限，但证明在此明确使用该条件。
关于 $P'\to Q'$ 的假设保证：对任意 $P$-模 $E$，映射
$E\to E\otimes_PQ$ 在 $J'$-幂挠子模上为双射；参见代数进阶引理
\ref{more-algebra-lemma-neighbourhood-equivalence}。
因此只需证明 $\omega(M)$ 是 $J'$-幂挠的。换言之，只需证明
$\omega(M)$ 在
$$
\text{Ext}^2_P(M, M \otimes_P IP)_g =
\text{Ext}^2_{P_g}(M_g, M_g \otimes_{P_g} IP_g)
$$
中对所有 $g\in J'$ 都变为零。由 $\omega(M)$ 的构造与基变换的相容性，
这确实成立，因为假设 $M_g$ 有提升（当然还须再次使用整串等价）。
% P11-E0056：冻结源将 However 误拼为 Howeover。
\end{proof}

\begin{lemma}
\label{examples-defos-lemma-lift-equivalence}
设 $A'\to A$ 是核幂零的 Noether 环满射，
$A\to B$ 是有限型平坦环同态。设 $\mathfrak b\subset B$ 是理想，
使得 $\Spec(B)\to\Spec(A)$ 在 $V(\mathfrak b)$ 的补集上为 syntomic。
则 $B$ 有到 $A'$ 的平坦提升，当且仅当其 $\mathfrak b$-进完备化
$B^\wedge$ 有到 $A'$ 的平坦提升。
\end{lemma}

\begin{proof}
选取 $A$-代数满射 $P=A[x_1,\ldots,x_n]\to B$。
令 $\mathfrak p\subset P$ 为 $\mathfrak b$ 的逆像。置
$P'=A'[x_1,\ldots,x_n]$，并以 $\mathfrak p'\subset P'$
表示 $\mathfrak p$ 的逆像（当然，这里的 $\mathfrak p$ 与
$\mathfrak p'$ 并不表示素理想）。以 $P^\wedge$ 与 $(P')^\wedge$
分别表示相应的完备化。

\medskip\noindent
假设 $A'\to B'$ 是 $A\to B$ 的平坦提升；换言之，
$A'\to B'$ 平坦，且有 $A$-代数同构 $B=B'\otimes_{A'}A$。
于是可选取提升给定满射 $P\to B$ 的 $A'$-代数同态 $P'\to B'$。
由 Nakayama 引理（代数引理 \ref{algebra-lemma-NAK}），
$B'$ 是 $P'$ 的商。因此可以赋予 $B'$ 一个在 $A'$ 上平坦的
$P'$-模结构，它提升作为在 $A$ 上平坦的 $P$-模的 $B$。
反之，若可把 $B$ 提升为在 $A'$ 上平坦的 $P'$-模 $M'$，
则再次应用 Nakayama 可知 $M'$ 是循环模 $M'\cong P'/J'$；
置 $B'=P'/J'$，便得到 $B$ 作为代数的平坦提升。

\medskip\noindent
置 $C=B^\wedge$ 与 $\mathfrak c=\mathfrak bC$。
假设 $A'\to C'$ 是 $A\to C$ 的平坦提升。则 $C'$ 关于
$\mathfrak c$ 的逆像 $\mathfrak c'$ 完备（代数引理
\ref{algebra-lemma-complete-modulo-nilpotent}）。选取提升
$A$-代数同态 $P\to C$ 的 $A'$-代数同态 $P'\to C'$。
这些同态经由完备化给出满射 $P^\wedge\to C$ 与
$(P')^\wedge\to C'$（对后一个再次使用 Nakayama 引理）。
特别地，可以赋予 $C'$ 一个在 $A'$ 上平坦的 $(P')^\wedge$-模结构，
它提升作为在 $A$ 上平坦的 $P^\wedge$-模的 $C$。
反之，若可把 $C$ 提升为在 $A'$ 上平坦的 $(P')^\wedge$-模 $N'$，
则再次应用 Nakayama 可知 $N'$ 是循环模
$N'\cong(P')^\wedge/\tilde J$；置
$C'=(P')^\wedge/\tilde J$，便得到 $C$ 作为代数的平坦提升。

\medskip\noindent
注意，$P'\to(P')^\wedge$ 是平坦环同态，并诱导同构
$P'/\mathfrak p'=(P')^\wedge/\mathfrak p'(P')^\wedge$。
因此，只要能证明对 $g\in\mathfrak p'$，$B_g$ 可提升为
在 $A'$ 上平坦的 $P'_g$-模，本引理便由引理
\ref{examples-defos-lemma-lift-equivalence-module} 推出。然而环同态
$A\to B_g$ 是 syntomic，故由环同态的光滑化命题
\ref{smoothing-proposition-lift-smooth}，它可提升为在 $A'$ 上平坦的代数
$B'$。由于 $A'\to P'_g$ 光滑，可以像此前一样把
$P_g\to B_g$ 提升为满射 $P'_g\to B'$，这正是所需结论。
\end{proof}

\noindent
记号。设 $A\to B$ 是环同态，$N$ 是 $B$-模。以
$\text{Exal}_A(B,N)$ 表示扩张
$$
0 \to N \to C \to B \to 0
$$
的同构类所成的集合；这里扩张属于 $A$-代数，且 $N$ 是 $C$
中的平方零理想。给定另一个这样的扩张
$0\to N\to C'\to B\to0$，所谓同构，是使得下图交换的
$A$-代数同构 $C\to C'$：
$$
\xymatrix{
0 \ar[r] &
N \ar[r] \ar[d]_{\text{id}} &
C \ar[r] \ar[d] &
B \ar[r] \ar[d]_{\text{id}} & 0 \\
0 \ar[r] &
N \ar[r] &
C' \ar[r] &
B \ar[r] & 0
}
$$
指派 $N\mapsto\text{Exal}_A(B,N)$ 是把积变成积的函子。
因此它是加性函子，而 $\text{Exal}_A(B,N)$ 具有自然的
$B$-模结构。事实上，由形变理论引理 \ref{defos-lemma-choices}，
有 $\text{Exal}_A(B,N)=\text{Ext}^1_B(\NL_{B/A},N)$。

\begin{lemma}
\label{examples-defos-lemma-first-order-completion}
设 $k$ 是域，$B$ 是有限型 $k$-代数。设 $J\subset B$ 是理想，
使得 $\Spec(B)\to\Spec(k)$ 在 $V(J)$ 的补集上光滑。
设 $N$ 是有限 $B$-模。则有典范双射
$$
\text{Exal}_k(B, N) \to \text{Exal}_k(B^\wedge, N^\wedge)
$$
。这里 $B^\wedge$ 与 $N^\wedge$ 是 $J$-进完备化。
\end{lemma}

\begin{proof}
该映射由完备化给出：给定 $\text{Exal}_k(B,N)$ 中的
$0\to N\to C\to B\to0$，把它送到 $C$ 关于 $J$ 的逆像的
完备化 $C^\wedge$。可与引理 \ref{examples-defos-lemma-completion} 的证明比较。

\medskip\noindent
由于 $k\to B$ 是有限呈示，复形 $\NL_{B/k}$ 可由复形
$N^{-1}\to N^0$ 表示，其中每个 $N^i$ 都是有限 $B$-模；
参见代数第 \ref{algebra-section-netherlander} 节，特别是代数引理
\ref{algebra-lemma-NL-homotopy}。由于 $B$ 是 Noether 环，
这意味着 $\NL_{B/k}$ 是伪凝聚的。对 $g\in J$，$k$-代数
$B_g$ 光滑，因而 $(\NL_{B/k})_g=\NL_{B_g/k}$ 拟同构于置于次数
$0$ 的有限投射 $B$-模。
% P11-E0058：冻结源此处的局部化复形自然应是有限投射 B_g-模，而非 B-模。
所以对任意 $B$-模 $N$ 及 $i\geq1$，有
$\text{Ext}^i_B(\NL_{B/k},N)_g=0$。由代数进阶引理
\ref{more-algebra-lemma-ext-annihilated-into} 可得
$$
\text{Ext}^1_B(\NL_{B/k}, N) \longrightarrow
\lim_n \text{Ext}^1_B(\NL_{B/k}, N/J^n N)
$$
对任意有限 $B$-模 $N$ 都是同构。

\medskip\noindent
映射的单射性。假设
$0\to N\to C\to B\to0$ 属于 $\text{Exal}_k(B,N)$，
并在 $\text{Exal}_k(B^\wedge,N^\wedge)$ 中映为零。
选取分裂 $C^\wedge=B^\wedge\oplus N^\wedge$。于是诱导映射
$C\to C^\wedge\to N^\wedge$ 对每个 $n$ 都给出映射
$C\to N/J^nN$。因此可见，我们的元素属于各映射
$$
\text{Ext}^1_B(\NL_{B/k}, N) \to
\text{Ext}^1_B(\NL_{B/k}, N/J^n N)
$$
的核，且这对所有 $n$ 均成立。由上一段可知，该元素为零。

\medskip\noindent
映射的满射性。设
$0\to N^\wedge\to C'\to B^\wedge\to0$ 是
$\text{Exal}_k(B^\wedge,N^\wedge)$ 中的元素。
沿 $B\to B^\wedge$ 拉回，得到 $\text{Exal}_k(B,N^\wedge)$
中的元素 $0\to N^\wedge\to C''\to B\to0$。有
$$
\text{Ext}^1_B(\NL_{B/k}, N^\wedge) =
\text{Ext}^1_B(\NL_{B/k}, N) \otimes_B B^\wedge =
\text{Ext}^1_B(\NL_{B/k}, N)
$$
。第一个等式由 $N^\wedge=N\otimes_BB^\wedge$（代数引理
\ref{algebra-lemma-completion-tensor}）以及代数进阶引理
\ref{more-algebra-lemma-pseudo-coherence-and-ext} 得到。
第二个等式成立，是因为 $\text{Ext}^1_B(\NL_{B/k},N)$
是 $J$-幂挠的（见上文），$B\to B^\wedge$ 平坦并诱导同构
$B/J\to B^\wedge/JB^\wedge$，以及代数进阶引理
\ref{more-algebra-lemma-neighbourhood-equivalence}。因此可找到
$C\in\text{Exal}_k(B,N)$，它在
$\text{Exal}_k(B,N^\wedge)$ 中映到 $C''$。于是
$$
0 \to N^\wedge \to C' \to B^\wedge \to 0
\quad\text{且}\quad
0 \to N^\wedge \to C^\wedge \to B^\wedge \to 0
$$
是 $\text{Exal}_k(B^\wedge,N^\wedge)$ 中的两个元素，
它们映到 $\text{Exal}_k(B,N^\wedge)$ 中的同一个元素。
取差，得到 $\text{Exal}_k(B^\wedge,N^\wedge)$ 中的元素
$0\to N^\wedge\to C'\to B^\wedge\to0$，其在
$\text{Exal}_k(B,N^\wedge)$ 中的像为零。这意味着存在
$$
\xymatrix{
0 \ar[r] &
N^\wedge \ar[r] &
C' \ar[r] &
B^\wedge \ar[r] & 0 \\
& & B \ar[u]^\sigma \ar[ru]
}
$$
。令 $J'\subset C'$ 为 $JB^\wedge\subset B^\wedge$ 的逆像。
为完成证明，只需注意：$\sigma$ 关于 $B$ 上的 $J$-进拓扑和
$C'$ 上的 $J'$-进拓扑连续，而且由代数引理
\ref{algebra-lemma-complete-modulo-nilpotent}，$C'$ 是 $J'$-进完备的
（这里还使用了 $C'$ 是 Noether 环；略去一个小细节）。
这就是说，$\sigma$ 经由完备化 $B^\wedge$ 分解，且
$C'=0$ 在 $\text{Exal}_k(B^\wedge,N^\wedge)$ 中成立。
\end{proof}

\begin{lemma}
\label{examples-defos-lemma-smooth-completion}
在例 \ref{examples-defos-example-rings} 中，设 $P$ 是 $k$-代数，
$J\subset P$ 是理想，并以 $P^\wedge$ 表示 $J$-进完备化。若
\begin{enumerate}
\item $k\to P$ 是有限型的；
\item $\Spec(P)\to\Spec(k)$ 在 $V(J)$ 的补集上光滑，
\end{enumerate}
则引理 \ref{examples-defos-lemma-completion} 中形变范畴之间的函子
$$
\Deformationcategory_P \longrightarrow \Deformationcategory_{P^\wedge}
$$
光滑，并在切空间上诱导同构。
\end{lemma}

\begin{proof}
由引理 \ref{examples-defos-lemma-rings-RS}，我们知道
$\Deformationcategory_P$ 与 $\Deformationcategory_{P^\wedge}$
都是形变范畴。因此只须检验该函子等同切空间并对应可提升性；
参见形式形变理论引理 \ref{formal-defos-lemma-easy-check-smooth}。
可提升性由引理 \ref{examples-defos-lemma-lift-equivalence} 证明，
切空间上的同构则是引理 \ref{examples-defos-lemma-first-order-completion}
中 $N=B$ 的特殊情形。
\end{proof}



\section{局部化的形变}
\label{examples-defos-section-compare-localization}

\noindent
本节比较由一个代数及其关于乘法子集的局部化所给出的形变问题。
我们先讨论“可提升性”。

\begin{lemma}
\label{examples-defos-lemma-lift-equivalence-localization}
设 $A'\to A$ 是核幂零的 Noether 环满射，
$A\to B$ 是有限型平坦环同态。设 $S\subset B$ 是乘法子集，
且只要 $\Spec(B)\to\Spec(A)$ 在 $\mathfrak q$ 处不是 syntomic，
就有 $S\cap\mathfrak q=\emptyset$。则 $B$ 有到 $A'$ 的平坦提升，
当且仅当 $S^{-1}B$ 有到 $A'$ 的平坦提升。
\end{lemma}

\begin{proof}
本证明与引理 \ref{examples-defos-lemma-lift-equivalence} 的证明相同，但更容易。
建议读者略过。选取 $A$-代数满射
$P=A[x_1,\ldots,x_n]\to B$。令 $S_P\subset P$ 为 $S$ 的逆像。
置 $P'=A'[x_1,\ldots,x_n]$，并以 $S_{P'}\subset P'$
表示 $S_P$ 的逆像。

\medskip\noindent
假设 $A'\to B'$ 是 $A\to B$ 的平坦提升；换言之，
$A'\to B'$ 平坦，且有 $A$-代数同构 $B=B'\otimes_{A'}A$。
于是可选取提升给定满射 $P\to B$ 的 $A'$-代数同态 $P'\to B'$。
由 Nakayama 引理（代数引理 \ref{algebra-lemma-NAK}），
$B'$ 是 $P'$ 的商。特别地，可以赋予 $B'$ 一个在 $A'$ 上平坦的
$P'$-模结构，它提升作为在 $A$ 上平坦的 $P$-模的 $B$。
反之，若可把 $B$ 提升为在 $A'$ 上平坦的 $P'$-模 $M'$，
则再次应用 Nakayama 可知 $M'$ 是循环模 $M'\cong P'/J'$；
置 $B'=P'/J'$，便得到 $B$ 作为代数的平坦提升。

\medskip\noindent
置 $C=S^{-1}B$。假设 $A'\to C'$ 是 $A\to C$ 的平坦提升。
$C'$ 中映到 $C$ 的可逆元的元素本身也可逆。选取提升
$A$-代数同态 $P\to C$ 的 $A'$-代数同态 $P'\to C'$。
由上述说明，这些同态经由局部化给出满射
$S_P^{-1}P\to C$ 与 $S_{P'}^{-1}P'\to C'$（对后一个使用
Nakayama 引理）。特别地，可以赋予 $C'$ 一个在 $A'$ 上平坦的
$S_{P'}^{-1}P'$-模结构，它提升作为在 $A$ 上平坦的
$S_P^{-1}P$-模的 $C$。反之，若可把 $C$ 提升为在 $A'$ 上平坦的
$S_{P'}^{-1}P'$-模 $N'$，则再次应用 Nakayama 可知 $N'$ 是循环模
$N'\cong S_{P'}^{-1}P'/\tilde J$；置
$C'=S_{P'}^{-1}P'/\tilde J$，便得到 $C$ 作为代数的平坦提升。

\medskip\noindent
概形态射的 syntomic 轨迹按定义是开集。令 $J_B\subset B$ 为
定义 $\Spec(B)$ 中 $\Spec(B)\to\Spec(A)$ 不是 syntomic 的点集的理想，
并以 $J_P\subset P$ 与 $J_{P'}\subset P'$ 表示相应理想。
注意，由引理中关于 $S$ 的假设，平坦环同态
$P'\to S_{P'}^{-1}P'$ 诱导同构
$P'/J_{P'}=S_{P'}^{-1}P'/J_{P'}S_{P'}^{-1}P'$；具体地说，
引理中的假设恰好是 $B/J_B=S^{-1}(B/J_B)$。
因此，只要能证明对 $g\in J_B$，$B_g$ 可提升为在 $A'$ 上平坦的
$P'_g$-模，本引理便由引理 \ref{examples-defos-lemma-lift-equivalence-module} 推出。
% P11-E0063：冻结源以 g∈J_B 开始，却随后写 P'_g；g 并非 P' 中已选定的元素。
然而环同态 $A\to B_g$ 是 syntomic，故由环同态的光滑化命题
\ref{smoothing-proposition-lift-smooth}，它可提升为在 $A'$ 上平坦的
代数 $B'$。由于 $A'\to P'_g$ 光滑，可以像此前一样把
$P_g\to B_g$ 提升为满射 $P'_g\to B'$，这正是所需结论。
\end{proof}

\begin{lemma}
\label{examples-defos-lemma-first-order-localization}
设 $k$ 是域，$B$ 是有限型 $k$-代数。设 $S\subset B$ 是乘法子集，
且只要 $\Spec(B)\to\Spec(k)$ 在 $\mathfrak q$ 处不光滑，
就有 $S\cap\mathfrak q=\emptyset$。设 $N$ 是有限 $B$-模。
则有典范双射
% P11-E0060：冻结源将 S 同时称为 multiplicative subset ideal，多出 ideal。
$$
\text{Exal}_k(B, N) \to \text{Exal}_k(S^{-1}B, S^{-1}N)
$$
\end{lemma}

\begin{proof}
本证明与引理 \ref{examples-defos-lemma-first-order-completion} 的证明相同，但更容易。
建议读者略过。该映射由局部化给出：给定
$\text{Exal}_k(B,N)$ 中的 $0\to N\to C\to B\to0$，
把它送到 $C$ 关于 $S$ 的逆像 $S_C\subset C$ 的局部化
$S_C^{-1}C$。可与引理 \ref{examples-defos-lemma-localization} 的证明比较。

\medskip\noindent
概形态射的光滑轨迹按定义是开集。令 $J\subset B$ 为定义
$\Spec(B)$ 中 $\Spec(B)\to\Spec(A)$ 不光滑的点集的理想。
% P11-E0061：本引理只定义了 k，此处冻结源却写 Spec(A)。
由于 $k\to B$ 是有限呈示，复形 $\NL_{B/k}$ 可由复形
$N^{-1}\to N^0$ 表示，其中每个 $N^i$ 都是有限 $B$-模；
参见代数第 \ref{algebra-section-netherlander} 节，特别是代数引理
\ref{algebra-lemma-NL-homotopy}。由于 $B$ 是 Noether 环，
这意味着 $\NL_{B/k}$ 是伪凝聚的。对 $g\in J$，$k$-代数
$B_g$ 光滑，因而 $(\NL_{B/k})_g=\NL_{B_g/k}$ 拟同构于置于次数
$0$ 的有限投射 $B$-模。
% P11-E0062：冻结源此处的局部化复形自然应是有限投射 B_g-模，而非 B-模。
所以对任意 $B$-模 $N$ 及 $i\geq1$，有
$\text{Ext}^i_B(\NL_{B/k},N)_g=0$。最后，有
$$
\text{Ext}^1_{S^{-1}B}(\NL_{S^{-1}B/k}, S^{-1}N) =
\text{Ext}^1_B(\NL_{B/k}, N) \otimes_B S^{-1}B =
\text{Ext}^1_B(\NL_{B/k}, N)
$$
。第一个等式由代数进阶引理
\ref{more-algebra-lemma-base-change-RHom} 与代数引理
\ref{algebra-lemma-localize-NL} 得到。第二个等式成立，是因为
$\text{Ext}^1_B(\NL_{B/k},N)$ 是 $J$-幂挠的，而 $S$ 中元素在
$J$-幂挠模上可逆地作用。结合引理 \ref{examples-defos-lemma-first-order-completion}
紧前方把 $\text{Exal}_A(B,N)$ 描述为
$\text{Ext}^1_B(\NL_{B/A},N)$ 的结论，证明完毕。
\end{proof}

\begin{lemma}
\label{examples-defos-lemma-smooth-localization}
在例 \ref{examples-defos-example-rings} 中，设 $P$ 是 $k$-代数，
$S\subset P$ 是乘法子集。若
\begin{enumerate}
\item $k\to P$ 是有限型的；
\item 对每个 $g\in S$，$\Spec(P)\to\Spec(k)$
在 $V(g)$ 的所有点处都光滑，
\end{enumerate}
则引理 \ref{examples-defos-lemma-localization} 中形变范畴之间的函子
$$
\Deformationcategory_P \longrightarrow \Deformationcategory_{S^{-1}P}
$$
光滑，并在切空间上诱导同构。
\end{lemma}

\begin{proof}
由引理 \ref{examples-defos-lemma-rings-RS}，我们知道
$\Deformationcategory_P$ 与 $\Deformationcategory_{S^{-1}P}$
都是形变范畴。因此只须检验该函子等同切空间并对应可提升性；
参见形式形变理论引理 \ref{formal-defos-lemma-easy-check-smooth}。
可提升性由引理 \ref{examples-defos-lemma-lift-equivalence-localization} 证明，
切空间上的同构则是引理 \ref{examples-defos-lemma-first-order-localization}
中 $N=B$ 的特殊情形。
\end{proof}



\section{Hensel 化的形变}
\label{examples-defos-section-compare-henselization}

\noindent
本节比较由一个代数及其完备化所给出的形变问题。
% P11-E0064：本节讨论 Hensel 化，冻结源却在此误写 completion。
我们先讨论“可提升性”。

\begin{lemma}
\label{examples-defos-lemma-lift-equivalence-henselization}
设 $A'\to A$ 是核幂零的 Noether 环满射，
$A\to B$ 是有限型平坦环同态。设 $\mathfrak b\subset B$ 是理想，
使得 $\Spec(B)\to\Spec(A)$ 在 $V(\mathfrak b)$ 的补集上为 syntomic。
设 $(B^h,\mathfrak b^h)$ 是对 $(B,\mathfrak b)$ 的 Hensel 化。
则 $B$ 有到 $A'$ 的平坦提升，当且仅当 $B^h$ 有到 $A'$ 的平坦提升。
\end{lemma}

\begin{proof}[第一种证明]
这个证明有所取巧。事实上，若 $B$ 有平坦提升 $B'$，
则取 Hensel 化 $(B')^h$，便得到 $B^h$ 的平坦提升
（可与引理 \ref{examples-defos-lemma-henselization} 的证明比较）。
反之，假设 $C'$ 是 $(B')^h$ 的一个在 $A'$ 上平坦的提升。
% P11-E0065：反向论证中冻结源写 (B')^h，但此处应为待提升的 B^h。
令 $\mathfrak c'\subset C'$ 为理想 $\mathfrak b^h$ 的逆像。
则 $C'$ 关于 $\mathfrak c'$ 的完备化 $(C')^\wedge$
是 $B^\wedge$ 的提升（细节从略）。因此由引理
\ref{examples-defos-lemma-lift-equivalence} 可知，$B$ 有平坦提升。
\end{proof}

\begin{proof}[第二种证明]
选取 $A$-代数满射 $P=A[x_1,\ldots,x_n]\to B$。
令 $\mathfrak p\subset P$ 为 $\mathfrak b$ 的逆像。置
$P'=A'[x_1,\ldots,x_n]$，并以 $\mathfrak p'\subset P'$
表示 $\mathfrak p$ 的逆像（当然，这里的 $\mathfrak p$ 与
$\mathfrak p'$ 并不表示素理想）。以 $P^h$ 与 $(P')^h$
分别表示相应的 Hensel 化。我们将使用 Hensel 化的函子性，
以及商的 Hensel 化就是 Hensel 化的相应商；参见代数进阶引理
\ref{more-algebra-lemma-irreducible-henselian-pair-connected} 与
\ref{more-algebra-lemma-henselization-integral}。

\medskip\noindent
假设 $A'\to B'$ 是 $A\to B$ 的平坦提升；换言之，
$A'\to B'$ 平坦，且有 $A$-代数同构 $B=B'\otimes_{A'}A$。
于是可选取提升给定满射 $P\to B$ 的 $A'$-代数同态 $P'\to B'$。
由 Nakayama 引理（代数引理 \ref{algebra-lemma-NAK}），
$B'$ 是 $P'$ 的商。特别地，可以赋予 $B'$ 一个在 $A'$ 上平坦的
$P'$-模结构，它提升作为在 $A$ 上平坦的 $P$-模的 $B$。
反之，若可把 $B$ 提升为在 $A'$ 上平坦的 $P'$-模 $M'$，
则再次应用 Nakayama 可知 $M'$ 是循环模 $M'\cong P'/J'$；
置 $B'=P'/J'$，便得到 $B$ 作为代数的平坦提升。

\medskip\noindent
置 $C=B^h$ 与 $\mathfrak c=\mathfrak bC$。假设 $A'\to C'$
是 $A\to C$ 的平坦提升。则 $C'$ 关于 $\mathfrak c$ 的逆像
$\mathfrak c'$ 是 Hensel 的（由代数进阶引理
\ref{more-algebra-lemma-henselian-henselian-pair} 以及
$C'\to C$ 的核幂零这一事实）。选取提升 $A$-代数同态
$P\to C$ 的 $A'$-代数同态 $P'\to C'$。这些同态经由 Hensel 化
给出满射 $P^h\to C$ 与 $(P')^h\to C'$（对后一个再次使用
Nakayama 引理）。特别地，可以赋予 $C'$ 一个在 $A'$ 上平坦的
$(P')^h$-模结构，它提升作为在 $A$ 上平坦的 $P^h$-模的 $C$。
反之，若可把 $C$ 提升为在 $A'$ 上平坦的 $(P')^h$-模 $N'$，
则再次应用 Nakayama 可知 $N'$ 是循环模
$N'\cong(P')^h/\tilde J$；置 $C'=(P')^h/\tilde J$，
便得到 $C$ 作为代数的平坦提升。

\medskip\noindent
注意，$P'\to(P')^h$ 是平坦环同态，并诱导同构
$P'/\mathfrak p'=(P')^h/\mathfrak p'(P')^h$（代数进阶引理
\ref{more-algebra-lemma-henselization-flat}）。因此，只要能证明对
$g\in\mathfrak p'$，$B_g$ 可提升为在 $A'$ 上平坦的 $P'_g$-模，
本引理便由引理 \ref{examples-defos-lemma-lift-equivalence-module} 推出。然而环同态
$A\to B_g$ 是 syntomic，故由环同态的光滑化命题
\ref{smoothing-proposition-lift-smooth}，它可提升为在 $A'$ 上平坦的
代数 $B'$。由于 $A'\to P'_g$ 光滑，可以像此前一样把
$P_g\to B_g$ 提升为满射 $P'_g\to B'$，这正是所需结论。
\end{proof}

\begin{lemma}
\label{examples-defos-lemma-first-order-henselization}
设 $k$ 是域，$B$ 是有限型 $k$-代数。设 $J\subset B$ 是理想，
使得 $\Spec(B)\to\Spec(k)$ 在 $V(J)$ 的补集上光滑。
设 $N$ 是有限 $B$-模。则有典范双射
$$
\text{Exal}_k(B, N) \to \text{Exal}_k(B^h, N^h)
$$
。这里 $(B^h,J^h)$ 是对 $(B,J)$ 的 Hensel 化，
且 $N^h=N\otimes_BB^h$。
\end{lemma}

\begin{proof}
本证明与引理 \ref{examples-defos-lemma-first-order-completion} 的证明相同，但更容易。
建议读者略过。该映射由 Hensel 化给出：给定
$\text{Exal}_k(B,N)$ 中的 $0\to N\to C\to B\to0$，
把它送到 $C$ 关于 $J$ 的逆像 $J_C\subset C$ 的 Hensel 化
$C^h$。可与引理 \ref{examples-defos-lemma-henselization} 的证明比较。

\medskip\noindent
由于 $k\to B$ 是有限呈示，复形 $\NL_{B/k}$ 可由复形
$N^{-1}\to N^0$ 表示，其中每个 $N^i$ 都是有限 $B$-模；
参见代数第 \ref{algebra-section-netherlander} 节，特别是代数引理
\ref{algebra-lemma-NL-homotopy}。由于 $B$ 是 Noether 环，
这意味着 $\NL_{B/k}$ 是伪凝聚的。对 $g\in J$，$k$-代数
$B_g$ 光滑，因而 $(\NL_{B/k})_g=\NL_{B_g/k}$ 拟同构于置于次数
$0$ 的有限投射 $B$-模。
% P11-E0066：冻结源此处的局部化复形自然应是有限投射 B_g-模，而非 B-模。
所以对任意 $B$-模 $N$ 及 $i\geq1$，有
$\text{Ext}^i_B(\NL_{B/k},N)_g=0$。最后，有
\begin{align*}
\text{Ext}^1_{B^h}(\NL_{B^h/k}, N^h)
& =
\text{Ext}^1_{B^h}(\NL_{B/k} \otimes_B B^h, N \otimes_B B^h) \\
& =
\text{Ext}^1_B(\NL_{B/k}, N) \otimes_B B^h \\
& =
\text{Ext}^1_B(\NL_{B/k}, N)
\end{align*}
第一个等式由代数进阶引理
\ref{more-algebra-lemma-henselization-NL}（更准确地说，是其关于
Hensel 对之 Hensel 化的类似结论）得到。第二个等式由代数进阶引理
\ref{more-algebra-lemma-base-change-RHom} 得到。第三个等式成立，
是因为 $\text{Ext}^1_B(\NL_{B/k},N)$ 是 $J$-幂挠的，映射
$B\to B^h$ 平坦并诱导同构 $B/J\to B^h/JB^h$（代数进阶引理
\ref{more-algebra-lemma-henselization-flat}），以及代数进阶引理
\ref{more-algebra-lemma-neighbourhood-equivalence}。结合引理
\ref{examples-defos-lemma-first-order-completion} 紧前方把
$\text{Exal}_A(B,N)$ 描述为 $\text{Ext}^1_B(\NL_{B/A},N)$ 的结论，
证明完毕。
\end{proof}

\begin{lemma}
\label{examples-defos-lemma-smooth-henselization}
在例 \ref{examples-defos-example-rings} 中，设 $P$ 是 $k$-代数，
$J\subset P$ 是理想。以 $(P^h,J^h)$ 表示对 $(P,J)$ 的 Hensel 化。若
\begin{enumerate}
\item $k\to P$ 是有限型的；
\item $\Spec(P)\to\Spec(k)$ 在 $V(J)$ 的补集上光滑，
\end{enumerate}
则引理 \ref{examples-defos-lemma-henselization} 中形变范畴之间的函子
$$
\Deformationcategory_P \longrightarrow \Deformationcategory_{P^h}
$$
光滑，并在切空间上诱导同构。
\end{lemma}

\begin{proof}
由引理 \ref{examples-defos-lemma-rings-RS}，我们知道
$\Deformationcategory_P$ 与 $\Deformationcategory_{P^h}$ 都是形变范畴。
因此只须检验该函子等同切空间并对应可提升性；参见形式形变理论引理
\ref{formal-defos-lemma-easy-check-smooth}。可提升性由引理
\ref{examples-defos-lemma-lift-equivalence-henselization} 证明，切空间上的同构则是
引理 \ref{examples-defos-lemma-first-order-henselization} 中 $N=B$ 的特殊情形。
\end{proof}



\section{孤立奇点的应用}
\label{examples-defos-section-isolated}

\noindent
我们应用上述讨论，研究只有有限多个奇点的有限型代数的形变理论。

\begin{lemma}
\label{examples-defos-lemma-isolated}
在例 \ref{examples-defos-example-rings} 中，设 $P$ 是 $k$-代数。
假设 $k\to P$ 是有限型的，且 $\Spec(P)\to\Spec(k)$
除在 $P$ 的极大理想 $\mathfrak m_1,\ldots,\mathfrak m_n$ 处外均光滑。
分别以 $P_{\mathfrak m_i}$、$P_{\mathfrak m_i}^h$、
$P_{\mathfrak m_i}^\wedge$ 表示局部环、Hensel 化和完备化。
则形变范畴的映射
$$
\Deformationcategory_P \to
\prod \Deformationcategory_{P_{\mathfrak m_i}} \to
\prod \Deformationcategory_{P_{\mathfrak m_i}^h} \to
\prod \Deformationcategory_{P_{\mathfrak m_i}^\wedge}
$$
光滑，并在其有限维切空间上诱导同构。
\end{lemma}

\begin{proof}
由引理 \ref{examples-defos-lemma-finite-type-rings-TI}，切空间是有限维的。
这些范畴之间的函子在引理 \ref{examples-defos-lemma-localization}、
\ref{examples-defos-lemma-henselization} 与 \ref{examples-defos-lemma-completion} 中构造
（我们略去一些形如“Hensel 化的完备化就是完备化”的验证）。

\medskip\noindent
置 $J=\mathfrak m_1\cap\ldots\cap\mathfrak m_n$，应用引理
\ref{examples-defos-lemma-smooth-completion} 可知
$\Deformationcategory_P \to \Deformationcategory_{P^\wedge}$
光滑，并在切空间上诱导同构；这里 $P^\wedge$ 是 $P$ 的
$J$-进完备化。然而，由于 $P^\wedge=\prod P_{\mathfrak m_i}^\wedge$，
可见映射 $\Deformationcategory_P\to
\prod \Deformationcategory_{P_{\mathfrak m_i}^\wedge}$
光滑，并在切空间上诱导同构。

\medskip\noindent
设 $(P^h,J^h)$ 是对 $(P,J)$ 的 Hensel 化。则
$P^h=\prod P_{\mathfrak m_i}^h$（考察幂等元并应用代数进阶引理
\ref{more-algebra-lemma-characterize-henselian-pair}）。
因此可应用引理 \ref{examples-defos-lemma-smooth-henselization}，
像完备化情形一样得到结论。

\medskip\noindent
为得到最后一种情形，只需分别对每个 $i$ 证明
$\Deformationcategory_{P_{\mathfrak m_i}} \to
\Deformationcategory_{P_{\mathfrak m_i}^\wedge}$
光滑，并在切空间上诱导同构。
% P11-E0067：冻结源的单数主语 map 后写 induce，缺少第三人称单数 s。
为此，可把 $P$ 替换为一个主局部化，使其唯一奇点是极大理想
$\mathfrak m$（对应于原来 $P$ 中的 $\mathfrak m_i$）。
然后对乘法子集 $S=P\setminus\mathfrak m$ 应用引理
\ref{examples-defos-lemma-smooth-localization} 即得结论。略去次要细节。
\end{proof}






\section{无阻碍形变问题}
\label{examples-defos-section-unobstructed}

\noindent
设 $p:\mathcal{F}\to\mathcal{C}_\Lambda$ 是以群胚为纤维的余纤维化范畴。
回忆，若 $p$ 光滑，就称 $\mathcal{F}$ {\it 光滑}或{\it 无阻碍}。
这意味着：给定 $\mathcal{C}_\Lambda$ 中的满射
$\varphi:A'\to A$ 以及 $x\in\Ob(\mathcal{F}(A))$，
存在 $\mathcal{F}$ 中的态射 $f:x'\to x$，满足
$p(f)=\varphi$。参见形式形变理论第
\ref{formal-defos-section-smooth} 节。本节给出若干具有几何意义的例子。

\begin{lemma}
\label{examples-defos-lemma-lci-unobstructed}
在例 \ref{examples-defos-example-rings} 中，设 $P$ 是 $k$ 上的局部完全交
（代数定义 \ref{algebra-definition-lci-field}）。
则 $\Deformationcategory_P$ 无阻碍。
\end{lemma}

\begin{proof}
设 $(A,Q)\to(k,P)$ 是 $\Deformationcategory_P$ 的对象。
由代数定义 \ref{algebra-definition-lci}，$A\to Q$ 是 syntomic 环同态。
因此，对 $\mathcal{C}_\Lambda$ 中任意满射 $A'\to A$，
由环同态的光滑化命题 \ref{smoothing-proposition-lift-smooth}，
存在提升 $A'\to A$ 的态射 $(A',Q')\to(A,Q)$。这就证明了引理。
\end{proof}

\begin{lemma}
\label{examples-defos-lemma-glueing-smooth}
在情形 \ref{examples-defos-situation-glueing} 中，若 $U_{12}\to\Spec(k)$ 光滑，
则态射
$$
\Deformationcategory_X
\longrightarrow
\Deformationcategory_{U_1} \times \Deformationcategory_{U_2} =
\Deformationcategory_{P_1} \times \Deformationcategory_{P_2}
$$
光滑。若此外 $U_1$ 是 $k$ 上的局部完全交，则
$$
\Deformationcategory_X
\longrightarrow
\Deformationcategory_{U_2} = \Deformationcategory_{P_2}
$$
光滑。
\end{lemma}

\begin{proof}
这些等号由引理 \ref{examples-defos-lemma-affine} 成立。按照形式形变理论第
\ref{formal-defos-section-smooth} 节，把 $\mathcal{C}_\Lambda$
看作 $\mathcal{C}_\Lambda$ 上的形变范畴。于是
$$
\Deformationcategory_{P_1} \times \Deformationcategory_{P_2} =
\Deformationcategory_{P_1}
\times_{\mathcal{C}_\Lambda}
\Deformationcategory_{P_2},
$$
参见形式形变理论注
\ref{formal-defos-remarks-cofibered-groupoids}
（\ref{formal-defos-item-product}）。利用引理 \ref{examples-defos-lemma-glueing}，
第一个断言是函子
$$
\Deformationcategory_{P_1}
\times_{\Deformationcategory_{P_{12}}}
\Deformationcategory_{P_2}
\longrightarrow
\Deformationcategory_{P_1}
\times_{\mathcal{C}_\Lambda}
\Deformationcategory_{P_2}
$$
光滑。只要能证明 $T\Deformationcategory_{P_{12}}=(0)$，
这就由形式形变理论引理
\ref{formal-defos-lemma-map-fibre-products-smooth} 得到。
由于 $P_{12}$ 在 $k$ 上光滑，该消失性由引理
\ref{examples-defos-lemma-smooth} 得到。对第二个断言，只需证明
$\Deformationcategory_{P_1} \to \mathcal{C}_\Lambda$
光滑；参见形式形变理论引理
\ref{formal-defos-lemma-smooth-properties}。换言之，须证明
$\Deformationcategory_{P_1}$ 无阻碍，这正是引理
\ref{examples-defos-lemma-lci-unobstructed}。
\end{proof}

\begin{lemma}
\label{examples-defos-lemma-curve-isolated}
在例 \ref{examples-defos-example-schemes} 中，设 $X$ 是 $k$ 上的概形。假设
\begin{enumerate}
\item $X$ 是分离的、在 $k$ 上有限型，且 $\dim(X)\leq1$；
\item $X\to\Spec(k)$ 除在闭点 $p_1,\ldots,p_n\in X$ 处外均光滑。
\end{enumerate}
分别以 $\mathcal{O}_{X,p_1}$、$\mathcal{O}_{X,p_1}^h$、
$\mathcal{O}_{X,p_1}^\wedge$ 表示局部环、Hensel 化和完备化。
% P11-E0068：冻结源在此三次写 p_1，而随后乘积及证明均按 p_i 取遍各奇点。
考虑形变范畴的映射
$$
\Deformationcategory_X
\longrightarrow
\prod \Deformationcategory_{\mathcal{O}_{X, p_i}}
\longrightarrow
\prod \Deformationcategory_{\mathcal{O}_{X, p_i}^h}
\longrightarrow
\prod \Deformationcategory_{\mathcal{O}_{X, p_i}^\wedge}
$$
第一个箭头光滑；第二、第三个箭头光滑，并在切空间上诱导同构。
\end{lemma}

\begin{proof}
选取仿射开集 $U_2\subset X$，使其包含 $p_1,\ldots,p_n$
以及 $X$ 每个不可约分支的泛点。这由代数簇引理
\ref{varieties-lemma-dim-1-quasi-projective} 与性质引理
\ref{properties-lemma-ample-finite-set-in-affine} 得以实现。
于是 $X\setminus U_2$ 是有限集，可以选取仿射开集
$U_1\subset X\setminus\{p_1,\ldots,p_n\}$，使得
$X=U_1\cup U_2$。置 $U_{12}=U_1\cap U_2$。
此时 $U_1$ 与 $U_{12}$ 都是 $k$ 上的光滑仿射概形。因此
$$
\Deformationcategory_X \longrightarrow \Deformationcategory_{U_2}
$$
由引理 \ref{examples-defos-lemma-glueing-smooth} 是光滑的。
应用引理 \ref{examples-defos-lemma-affine} 与 \ref{examples-defos-lemma-isolated} 即得结论。
\end{proof}

\begin{lemma}
\label{examples-defos-lemma-curve-isolated-lci}
在例 \ref{examples-defos-example-schemes} 中，设 $X$ 是 $k$ 上的概形。假设
\begin{enumerate}
\item $X$ 是分离的、在 $k$ 上有限型，且 $\dim(X)\leq1$；
\item $X$ 是 $k$ 上的局部完全交；
\item $X\to\Spec(k)$ 除在有限多个点处外均光滑。
\end{enumerate}
则 $\Deformationcategory_X$ 无阻碍。
\end{lemma}

\begin{proof}
设 $p_1,\ldots,p_n\in X$ 是 $X\to\Spec(k)$ 不光滑的点。
选取仿射开集 $U_2\subset X$，使其包含 $p_1,\ldots,p_n$
以及 $X$ 每个不可约分支的泛点。这由代数簇引理
\ref{varieties-lemma-dim-1-quasi-projective} 与性质引理
\ref{properties-lemma-ample-finite-set-in-affine} 得以实现。
于是 $X\setminus U_2$ 是有限集，可以选取仿射开集
$U_1\subset X\setminus\{p_1,\ldots,p_n\}$，使得
$X=U_1\cup U_2$。置 $U_{12}=U_1\cap U_2$。
此时 $U_1$ 与 $U_{12}$ 都是 $k$ 上的光滑仿射概形。因此
$$
\Deformationcategory_X \longrightarrow \Deformationcategory_{U_2}
$$
由引理 \ref{examples-defos-lemma-glueing-smooth} 是光滑的。
应用引理 \ref{examples-defos-lemma-affine} 与 \ref{examples-defos-lemma-lci-unobstructed} 即得结论。
\end{proof}




\section{光滑化}
\label{examples-defos-section-smoothing}

\noindent
设给定域 $k$ 上的有限型概形或代数空间 $X$。寻找有限型平坦态射
$Y\to\Spec(k[[t]])$ 往往很有用，其中一般纤维光滑，特殊纤维同构于
$X$。这样的对象称为 $X$ 的一个光滑化。本节将为具有孤立局部完全交
奇点的 $1$ 维分离 $X$ 找到光滑化。

\begin{lemma}
\label{examples-defos-lemma-criterion-smoothing}
设 $k$ 是域。置 $S=\Spec(k[[t]])$ 与
$S_n=\Spec(k[t]/(t^n))$。设 $Y\to S$ 是固有平坦概形态射，
其特殊纤维 $X$ 是 Cohen--Macaulay 且为 $d$ 维等维的。
记 $X_n=Y\times_SS_n$。若对某个 $n\geq1$，
$\Omega_{X_n/S_n}$ 的第 $d$ 个 Fitting 理想包含 $t^{n-1}$，
则 $Y\to S$ 的一般纤维光滑。
\end{lemma}

\begin{proof}
由态射进阶引理
\ref{more-morphisms-lemma-flat-finite-presentation-CM-open}
可知 $Y\to S$ 是 Cohen--Macaulay 态射。由态射引理
\ref{morphisms-lemma-flat-finite-presentation-CM-fibres-relative-dimension}
可知 $Y\to S$ 的相对维数为 $d$。由除子引理
\ref{divisors-lemma-d-fitting-ideal-omega-smooth}，
$\Omega_{Y/S}$ 的第 $d$ 个 Fitting 理想
$\mathcal{I}\subset\mathcal{O}_Y$ 定义了态射 $Y\to S$ 的奇异轨迹。
换言之，$V(\mathcal{I})\subset Y$ 是 $Y\to S$ 不光滑的点所成的闭子集。
由除子引理 \ref{divisors-lemma-base-change-and-fitting-ideal-omega}，
该 Fitting 理想的形成与基变换交换。根据假设，$t^{n-1}$ 是
$\mathcal{I}+t^n\mathcal{O}_Y$ 的截面。因此，对每个
$x\in X=V(t)\subset Y$，有 $t^{n-1}\in\mathcal{I}_x$；
这里 $\mathcal{I}_x$ 是在 $x$ 处的茎。这意味着在 $Y$ 中 $X$
的某个开邻域内有 $V(\mathcal{I})\subset V(t)$。由于 $Y\to S$
固有，可得所需的 $V(\mathcal{I})\subset V(t)$。
\end{proof}

\begin{lemma}
\label{examples-defos-lemma-jouanolou-type-thing}
设 $k$ 是域，$1\leq c\leq n$ 是整数。设
$f_1,\ldots,f_c\in k[x_1,\ldots x_n]$，并令
$a_{ij}$（$0\leq i\leq n$，$1\leq j\leq c$）为变量。考虑
% P11-E0069：冻结源的多项式环变量表在 \ldots 与 x_n 之间缺逗号。
$$
g_j = f_j + a_{0j} + a_{1j}x_1 + \ldots + a_{nj}x_n \in
k[a_{ij}][x_1, \ldots, x_n]
$$
以 $Y\subset\mathbf{A}^{n+c(n+1)}_k$ 表示由
$g_1,\ldots,g_c$ 定义的闭子概形。以
$\pi:Y\to\mathbf{A}^{c(n+1)}_k$ 表示到以 $a_{ij}$ 为变量的
仿射空间的投影。则 $\mathbf{A}^{c(n+1)}_k$ 中存在非空 Zariski
开子集，使得 $\pi$ 在其上光滑。
% P11-E0070：冻结源写 “a nonempty Zariski open of”，缺少 subset 等中心词。
\end{lemma}

\begin{proof}
回忆，$\pi$ 光滑的点所成的集合是开集，故其补集，即奇异轨迹，
是闭集。由 Chevalley 定理（采用态射引理
\ref{morphisms-lemma-chevalley} 的形式），奇异轨迹的像是可构造集。
因此，只要 $\mathbf{A}^{c(n+1)}_k$ 的泛点不在奇异轨迹的像中，
本引理便成立（例如可用拓扑引理
\ref{topology-lemma-generic-point-in-constructible}）。所以须证明：
不存在映到 $\mathbf{A}^{c(n+1)}_k$ 泛点的点 $y\in Y$，
使得 $\pi$ 在该点不光滑。考虑偏导数矩阵
$$
(\frac{\partial g_j}{\partial x_i}) =
(\frac{\partial f_j}{\partial x_i} + a_{ij})
$$
该矩阵在 $\kappa(y)$ 中的像必有秩 $<c$，否则 $\pi$ 将在 $y$
处光滑；参见环同态的光滑化第
\ref{smoothing-section-singular-ideal} 节的讨论。因此可找到不全为零的
$\lambda_1,\ldots,\lambda_c\in\kappa(y)$，使向量
$(\lambda_1,\ldots,\lambda_c)$ 位于该矩阵的核中。重新编号后，
可假设 $\lambda_1\not = 0$。除以 $\lambda_1$ 后，可假设该向量形如
$(1,\lambda_2,\ldots,\lambda_c)$。于是得到
% P11-E0071：冻结源下一公式首项含自由指标 j 且微分变量误为 x_1。
$$
a_{i1} = -
\frac{\partial f_j}{\partial x_1} -
\sum\nolimits_{j = 2, \ldots, c} \lambda_j(\frac{\partial f_j}{\partial x_i} + a_{ij})
$$
在 $\kappa(y)$ 中对 $i=1,\ldots,n$ 成立。此外，由于 $y\in Y$，还有
$$
a_{0j} = -f_j - a_{1j}x_1 - \ldots - a_{nj}x_n
$$
在 $\kappa(y)$ 中成立。这意味着 $\kappa(y)$ 中由 $a_{ij}$ 生成的
子域，包含于 $\kappa(y)$ 中由 $x_1,\ldots,x_n,\lambda_2,\ldots,\lambda_c$ 的像，
以及除 $a_{i1}$ 与 $a_{0j}$ 外的 $a_{ij}$ 所生成的子域。
计数可知，后者的超越次数至多为 $c(n+1)-1$。因此 $y$
不可能映到泛点，正如所需。
\end{proof}

\begin{lemma}
\label{examples-defos-lemma-smoothing-affine-lci}
设 $k$ 是域，$A$ 是 $k$ 上的全局完全交。存在有限型平坦环同态
$k[[t]]\to B$，满足 $B/tB\cong A$，并使 $B[1/t]$
在 $k((t))$ 上光滑。
\end{lemma}

\begin{proof}
按照代数定义 \ref{algebra-definition-lci-field}，写成
$A=k[x_1,\ldots,x_n]/(f_1,\ldots,f_c)$。我们将选取
$a_{ij}\in(t)\subset k[[t]]$，并置
$$
g_j = f_j + a_{0j} + a_{1j}x_1 + \ldots + a_{nj}x_n \in
k[[t]][x_1, \ldots, x_n]
$$
。随后取 $B=k[[t]][x_1,\ldots,x_n]/(g_1,\ldots,g_c)$。
我们断言 $k[[t]]\to B$ 在每个位于 $(t)$ 上方的素理想处都平坦。
事实上，元素 $f_1,\ldots,f_c$ 对 $k[x_1,\ldots,x_n]$ 中任何包含
$f_1,\ldots,f_c$ 的素理想 $\mathfrak p$，都在相应局部环中
构成正则序列（代数引理 \ref{algebra-lemma-lci}）。因此
$g_1,\ldots,g_c$ 局部地是正则序列的提升，可以应用代数引理
\ref{algebra-lemma-grothendieck-regular-sequence}。
在位于 $(0)\subset k[[t]]$ 上方的素理想处，平坦性自动成立，
因为 $k((t))=k[[t]]_{(0)}$ 是域。因此 $B$ 在 $k[[t]]$ 上平坦。

\medskip\noindent
只须再证明：适当选取 $a_{ij}$ 后，一般纤维 $B_{(0)}$
在 $k((t))$ 上光滑。为此，须证明可以选取 $a_{ij}$，使诱导态射
$$
(a_{ij}) : \Spec(k[[t]]) \longrightarrow \mathbf{A}^{c(n + 1)}_k
$$
映入引理 \ref{examples-defos-lemma-jouanolou-type-thing} 中的非空 Zariski 开集。
这是显然的，因为不存在关于 $a_{ij}$ 的非零多项式在
$(t)^{\oplus c(n+1)}$ 上恒为零（留作读者练习）。
\end{proof}

\begin{lemma}
\label{examples-defos-lemma-smoothing-artinian-lci}
设 $k$ 是域。设 $A$ 是有限维 $k$-代数，并且是 $k$ 上的局部完全交。
则存在有限平坦 $k[[t]]$-代数 $B$，满足 $B/tB\cong A$，
且 $B[1/t]$ 在 $k((t))$ 上为 \'etale。
\end{lemma}

\begin{proof}
由于 $A$ 是 Artin 环（代数引理
\ref{algebra-lemma-finite-dimensional-algebra}），可把 $A$ 写成局部
Artin 环的积（代数引理 \ref{algebra-lemma-artinian-finite-length}）。
因此只须证明 $A$ 为局部环的情形（这里使用了取主局部化保持
局部完全交性质；参见代数引理 \ref{algebra-lemma-localize-lci}）。
在这种情形下，$A$ 是全局完全交。考虑引理
\ref{examples-defos-lemma-smoothing-affine-lci} 中构造的代数 $B$。则
$k[[t]]\to B$ 在 $B$ 中唯一位于 $(t)$ 上方的素理想处拟有限
（代数定义 \ref{algebra-definition-quasi-finite}）。注意，$k[[t]]$
是 Hensel 局部环（代数引理 \ref{algebra-lemma-complete-henselian}）。
因此 $B=B'\times C$，其中 $B'$ 在 $k[[t]]$ 上有限，而 $C$
没有位于 $(t)$ 上方的素理想；参见代数引理
\ref{algebra-lemma-characterize-henselian}。于是 $B'$ 正是所求的环
（回忆，\'etale 等价于相对维数为 $0$ 的光滑）。
\end{proof}

\begin{lemma}
\label{examples-defos-lemma-smoothing-at-lci-point}
设 $k$ 是域，$A$ 是 $k$-代数。假设
\begin{enumerate}
\item $A$ 是本质有限型 $k$-局部环；
\item $A$ 是 $k$ 上的完全交
（代数定义 \ref{algebra-definition-lci-local-ring}）。
\end{enumerate}
置 $d=\dim(A)+\text{trdeg}_k(\kappa)$，其中 $\kappa$ 是 $A$
的剩余域。则存在整数 $n$ 以及本质有限型平坦环同态
$k[[t]]\to B$，满足 $B/tB\cong A$，并使 $t^n$ 属于
$\Omega_{B/k[[t]]}$ 的第 $d$ 个 Fitting 理想。
% P11-E0072：冻结源 “a flat, essentially of finite type ring map” 的修饰语次序不合英语语法。
\end{lemma}

\begin{proof}
由代数引理 \ref{algebra-lemma-lci-local}，可把 $A$ 写成 $k$
上的全局完全交 $P$ 在某个素理想 $\mathfrak p$ 处的局部化。
由代数引理 \ref{algebra-lemma-dimension-at-a-point-finite-type-field}，
注意到 $\dim(P)=d$。由引理 \ref{examples-defos-lemma-smoothing-affine-lci}，
可找到有限型平坦环同态 $k[[t]]\to Q$，使得 $P\cong Q/tQ$，
并使 $k((t))\to Q[1/t]$ 光滑。由该引理中 $Q$ 的构造可知，
$k[[t]]\to Q$ 是相对维数为 $d$ 的相对全局完全交；或者，
代数引理 \ref{algebra-lemma-syntomic} 表明 $Q$ 或 $Q$
的某个适当主局部化是这样的全局完全交。因此由除子引理
\ref{divisors-lemma-d-fitting-ideal-omega-smooth}，
$\Omega_{Q/k[[t]]}$ 的第 $d$ 个 Fitting 理想 $I\subset Q$
定义了 $\Spec(Q)\to\Spec(k[[t]])$ 的奇异轨迹。
于是对某个 $n$ 有 $t^n\in I$。令 $\mathfrak q\subset Q$
为 $\mathfrak p$ 的逆像，置 $B=Q_\mathfrak q$。引理得证。
\end{proof}

\begin{lemma}
\label{examples-defos-lemma-smoothing-proper-curve-isolated-lci}
设 $X$ 是域 $k$ 上的概形。假设
\begin{enumerate}
\item $X$ 在 $k$ 上固有；
\item $X$ 是 $k$ 上的局部完全交；
\item $X$ 的维数 $\leq1$；
\item $X\to\Spec(k)$ 除在有限多个点处外均光滑。
\end{enumerate}
则存在平坦射影态射 $Y\to\Spec(k[[t]])$，其一般纤维光滑，
特殊纤维同构于 $X$。
\end{lemma}

\begin{proof}
注意，$X$ 是 Cohen--Macaulay 的；参见代数引理
\ref{algebra-lemma-lci-CM}。因此 $X=X'\amalg X''$，其中
$\dim(X')=0$，而 $X''$ 是 $1$ 维等维的；参见态射引理
\ref{morphisms-lemma-flat-finite-presentation-CM-fibres-relative-dimension}。
由于 $X'$ 在 $k$ 上有限（代数簇引理
\ref{varieties-lemma-algebraic-scheme-dim-0}），由引理
\ref{examples-defos-lemma-smoothing-artinian-lci} 可找到
$Y'\to\Spec(k[[t]])$，其特殊纤维为 $X'$，一般纤维光滑。
因此只须对 $X''$ 证明本引理。以 $X''$ 替换 $X$ 后，
$X$ 是 Cohen--Macaulay 且为 $1$ 维等维的。

\medskip\noindent
我们将对情形 $\Lambda=k\to k$ 使用形变理论。令
$p_1,\ldots,p_r\in X$ 为 $X$ 的闭奇点，即 $X\to\Spec(k)$
不光滑的点。对每个 $i$，选取整数 $n_i$ 和本质有限型平坦环同态
$$
k[[t]] \longrightarrow B_i
$$
，满足 $B_i/tB_i\cong\mathcal{O}_{X,p_i}$，并使 $t^{n_i}$ 属于
$\Omega_{B_i/k[[t]]}$ 的第 $1$ 个 Fitting 理想。由引理
\ref{examples-defos-lemma-smoothing-at-lci-point}，这可以做到。注意，系统
$(B_i/t^nB_i)$ 定义了 $\Deformationcategory_{\mathcal{O}_{X,p_i}}$
在 $k[[t]]$ 上的形式对象。由引理 \ref{examples-defos-lemma-curve-isolated}，映射
$$
\Deformationcategory_X
\longrightarrow
\prod\nolimits_{i = 1, \ldots, r} \Deformationcategory_{\mathcal{O}_{X, p_i}}
$$
是形变范畴之间的光滑映射。因此由形式形变理论引理
\ref{formal-defos-lemma-smooth-morphism-essentially-surjective}
，$\Deformationcategory_X$ 中存在形式对象 $(X_n)$，
经上述箭头映到形式对象 $\prod_i(B_i/t^n)$。由空间态射进阶引理
\ref{spaces-more-morphisms-lemma-formal-algebraic-space-proper-reldim-1}
，存在 $k[[t]]$ 上的射影概形 $Y$ 以及相容同构
$Y\times_{\Spec(k[[t]])}\Spec(k[t]/(t^n))\cong X_n$。
由态射进阶引理
\ref{more-morphisms-lemma-check-flatness-on-infinitesimal-nbhds}
可知 $Y\to\Spec(k[[t]])$ 平坦。由于 $X$ 是 Cohen--Macaulay
且为 $1$ 维等维的，可应用引理 \ref{examples-defos-lemma-criterion-smoothing}
检验 $Y$ 的一般纤维光滑\footnote{警告：一般而言，$Y$ 在点
$p_i$ 处的局部环同构于 $B_i$ 这一说法{\bf 不}成立。
我们只知道两边都模去 $t^n$ 后该说法成立！}。
选取严格大于上述整数 $n_i$ 之最大值的 $n$。若能证明
$t^{n-1}$ 属于 $\Omega_{X_n/S_n}$ 的第一个 Fitting 理想，
其中 $S_n=\Spec(k[t]/(t^n))$，则证明完成。为此，只须在 $X_n$
每个闭点 $p$ 处的局部环中证明这一点。然而，若 $p$ 对应于
$X\to\Spec(k)$ 的光滑点，则 $\Omega_{X_n/S_n,p}$ 是秩为 $1$
的自由模，第一个 Fitting 理想等于该局部环。若对某个 $i$
有 $p=p_i$，则
$$
\Omega_{X_n/S_n, p_i} =
\Omega_{(B_i/t^nB_i)/(k[t]/(t^n))} =
\Omega_{B_i/k[[t]]}/t^n\Omega_{B_i/k[[t]]}
$$
由于取 Fitting 理想与基变换交换（我们已经使用过这一点；
在这里的代数情形中，它由代数进阶引理
\ref{more-algebra-lemma-fitting-ideal-basics} 得到），并且
$n-1\geq n_i$，可知 $t^{n-1}$ 属于该模作为
$B_i/t^nB_i$-模的 Fitting 理想，正如所需。
% P11-E0073：冻结源将 “we already used” 误写为 “with already used”。
\end{proof}

\begin{lemma}
\label{examples-defos-lemma-smoothing-curve-isolated-lci}
设 $k$ 是域，$X$ 是 $k$ 上的概形。假设
\begin{enumerate}
\item $X$ 是分离的、在 $k$ 上有限型，且 $\dim(X)\leq1$；
\item $X$ 是 $k$ 上的局部完全交；
\item $X\to\Spec(k)$ 除在有限多个点处外均光滑。
\end{enumerate}
则存在有限型平坦分离态射 $Y\to\Spec(k[[t]])$，其一般纤维光滑，
特殊纤维同构于 $X$。
\end{lemma}

\begin{proof}
若 $X$ 既约，则可像代数簇引理
\ref{varieties-lemma-reduced-dim-1-projective-completion} 中那样选取嵌入
$X\subset\overline{X}$。写成
$X=\overline{X}\setminus\{x_1,\ldots,x_n\}$，可见
$\mathcal{O}_{\overline{X},x_i}$ 是离散赋值环，因而特别是局部完全交
（代数定义 \ref{algebra-definition-lci-local-ring}）。因此
$\overline{X}$ 是 $k$ 上的局部完全交：在开集 $X$ 上这一点成立，
而在点 $x_i$ 处由代数引理 \ref{algebra-lemma-lci-local} 成立。
故可应用引理 \ref{examples-defos-lemma-smoothing-proper-curve-isolated-lci}，找到
射影平坦态射 $\overline{Y}\to\Spec(k[[t]])$，其一般纤维光滑，
特殊纤维为 $\overline{X}$。随后从 $\overline{Y}$ 中移去
$x_1,\ldots,x_n$，得到 $Y$。

\medskip\noindent
在一般情形下，写成 $X=X'\amalg X''$，其中
$\dim(X')=0$，而 $X''$ 是 $1$ 维等维的。
% P11-E0074：冻结源在 where 后多写了 with。
此时 $X''$ 既约，第一段适用于它。另一方面，可以像引理
\ref{examples-defos-lemma-smoothing-proper-curve-isolated-lci} 的证明中那样处理
$X'$。略去一些细节。
\end{proof}
